% !TeX root = ../main.tex
\section{Constitutive models and the entropy inequality}
\label{sec:multicomponent_solid_skeletons}
\label{sec:multicomponent_solid_mixture_rules}
\label{sec:constitutive_models_entropy_inequality}

Section~\ref{sec:virtual_power_derivation} derived the governing equations from
a minimal set of independent kinematic fields and process rates, with the
constraints supplying the remaining phase and component relations.  That
derivation was general across solid and fluid phases, but it was not
constitutively neutral: choices needed to define the modeled physics were
already imposed, most notably the electroquasistatic specialization in place of
full electrodynamics.  We now make the additional constitutive choices required
for porous media, including the shared solid-skeleton motion, the thermal and
free-energy state descriptions, and the admissible stress, transport, reaction,
and phase-transformation relations identified through the entropy inequality.

The phase fields introduced in \sref{virtual_power_derivation} allow every
phase to carry component mass fractions.  The solid phases retain that full compositional
structure.  Each \(s\in\mc S\) is a load-bearing multicomponent solid phase with
mass fractions \(\eta_s^\alpha\), while all solid phases share the common
skeleton trajectory
\begin{equation}
	\mbf x=\bs\chi(\mbf X,t),
	\qquad
	\mbf F=\nabla_{\mbf X}\bs\chi,
	\qquad
	J=\det\mbf F .
	\label{eq:common_solid_skeleton_motion}
\end{equation}
Adsorbed species, host minerals, hydrate cages, bound water, and hydrate guest
molecules share the skeleton velocity.  Their mechanical effect enters through solid composition,
stress-free swelling, thermal expansion, elastic moduli, and plastic flow
properties.

The thermal specialization distinguishes an equilibrated mobile-fluid
subsystem from the solid skeleton,
%
\begin{equation}
	\theta_\xi
	=
	\begin{cases}
		\theta_{\mc F}, & \xi=f\in\mc F,\\
		\theta_{\mc S}, & \xi=s\in\mc S .
	\end{cases}
	\label{eq:MC_two_temperature_specialization}
\end{equation}
%
Thus the fluid phases share \(\theta_{\mc F}\), whereas the load-bearing solid
phases share \(\theta_{\mc S}\).  This two-temperature
local-thermal-nonequilibrium model permits the fluid and solid subsystems to
have different temperatures and resolves heat exchange between them; either
subsystem may be hotter.
Multiphase geothermal models use the same common-fluid and rock-temperature grouping
\cite{faraz2021boiling}, while constituent mixture theory permits the more
general phase-temperature description
\cite{drumheller2000,hanyga2004thermal}.  Local thermal equilibrium is recovered
by imposing \(\theta_{\mc F}=\theta_{\mc S}\).

The constitutive exploitation is written in a single observer frame attached to
the solid skeleton.  Define the skeleton velocity and, for every spatial
thermodynamic state field \(\Gamma\), the ordinary-dot notation by
%
\begin{equation}
	\mbf v_{\mc S}
	=
	\left.
	\frac{\partial\bs\chi}{\partial t}
	\right\vert_{\mbf X},
	\qquad
	\dot\Gamma
	\equiv
	\frac{D_{\mc S}\Gamma}{Dt}
	=
	\frac{\partial\Gamma}{\partial t}
	+
	\mbf v_{\mc S}\cdot\nabla_{\mbf x}\Gamma .
	\label{eq:MC_skeleton_rate_convention}
\end{equation}
%
This convention applies in particular to
\(\theta_{\mc F}\), \(\theta_{\mc S}\), every \(\phi_\xi\),
\(\bar\rho_\xi\), \(\eta_\xi^\alpha\), and the phase thermodynamic
potentials, as well as the solid kinematic and internal-variable fields
\(\mbf A_s\), \(\bar{\mbf F}_s\), and their determinants.  The heavy dot marks
rates written following their own phase, including the component-insertion
rates \(\phasedot{\mc C}_\xi^\alpha\) and
\(\phasedot{c}_\xi^\alpha\).  Because all solid phases share the skeleton
trajectory, \(D_s/Dt=D_{\mc S}/Dt\), whereas mobile phases retain their phase
velocities.  For any phase-attributed spatial field
\(\Gamma_\xi\), the phase-following and skeleton-following rates are related
exactly by
%
\begin{equation}
	\phasedot{\Gamma}_\xi
	=
	\dot\Gamma_\xi
	+
	\left(
	\mbf v_\xi-\mbf v_{\mc S}
	\right)\cdot\nabla_{\mbf x}\Gamma_\xi .
	\label{eq:MC_phase_skeleton_transport_identity}
\end{equation}
%
Thus \(\mbf v_\xi-\mbf v_{\mc S}\) and the associated mass, momentum, energy,
and entropy fluxes represent relative phase transport explicitly.
This solid-observer convention is used in thermodynamic formulations of
multicomponent multiphase poroelastic flow
\cite[eq.~(2)]{seguin2019multicomponent} and is consistent with the
solid-observed saturation rate used by Hassanizadeh and Gray
\cite[eq.~(35)]{hassanizadeh1993thermodynamic}.

The pointwise volume-fraction constraint therefore has the solid-material rate
%
\begin{equation}
	\dot\phi_\xi
	\equiv
	\frac{D_{\mc S}\phi_\xi}{Dt}
	=
	\frac{\partial\phi_\xi}{\partial t}
	+
	\mbf v_{\mc S}\cdot\nabla_{\mbf x}\phi_\xi,
	\qquad
	\sum_\xi\dot\phi_\xi=0 .
	\label{eq:MC_skeleton_volume_fraction_rates}
\end{equation}
%
Individual \(\dot\phi_\xi\) may be nonzero while their sum vanishes.  The
mechanism-progress rate \(\dot r_{(m)}\) retains the ordinary dot as the Onsager
process variable shared by all participating phases.  The electrostatic
potential is a single Eulerian field, and the electrical power retains its
explicit field-rate and mechanical-transport terms.  The explicitly written
\(D_\xi\tau/Dt\) later projects the shared transfer multiplier along phase
trajectory \(\xi\) in the component-insertion compatibility relation.

\subsection{Solid kinematics and state variables}
\label{sec:MC_solid_kinematics}

Extending Drumheller's distension/true-deformation split for reacting mixtures
\cite{drumheller2000} and its porous-media form in Foster et
al.~\cite{foster2026reactingporous}, the solid kinematic split is
phase-specific but uses the common skeleton deformation,
%
\begin{equation}
	\mbf F=\mbf A_s\bar{\mbf F}_s,
	\qquad
	J=a_s\bar J_s,
	\qquad
	a_s=\det\mbf A_s,
	\qquad
	\bar J_s=\det\bar{\mbf F}_s,
	\qquad
	s\in\mc S .
	\label{eq:MC_solid_distension_decomposition}
\end{equation}
%
The distension gradient tensor \(\mbf A_s\) stores pore-space allocation, packing,
stress-free swelling, and reaction-induced accommodation between the common
skeleton motion and the true solid deformation \(\bar{\mbf F}_s\). Extending
Drumheller's elastic--plastic treatment of true deformation and distension
\cite{drumheller2000}, the elastic, plastic, and stress-free factors are
%
\begin{subequations}
	\label{eq:MC_solid_kinematic_splits}
	\begin{align}
		\mbf A_s
		 & =
		\mbf A_s^e\mbf A_s^p\mbf A_s^0,
		\label{eq:MC_solid_distension_split} \\
		\bar{\mbf F}_s
		 & =
		\bar{\mbf F}_s^e
		\bar{\mbf F}_s^p
		\bar{\mbf F}_s^0 .
		\label{eq:MC_solid_true_split}
	\end{align}
\end{subequations}
%
Here \(\mbf A_s^e\) and \(\bar{\mbf F}_s^e\) are elastic factors,
\(\mbf A_s^p\) and \(\bar{\mbf F}_s^p\) are plastic internal variables, and
\(\mbf A_s^0\) and \(\bar{\mbf F}_s^0\) are reversible stress-free maps that can
include temperature and solid-composition effects,
%
\begin{equation}
	\mbf A_s^0
	=
	\mbf A_s^0(\theta_{\mc S},\{\eta_s^\alpha\}),
	\qquad
	\bar{\mbf F}_s^0
	=
	\bar{\mbf F}_s^0(\theta_{\mc S},\{\eta_s^\alpha\}).
	\label{eq:MC_solid_stress_free_maps}
\end{equation}
%
The stress-free maps represent adsorption swelling, desorption shrinkage,
hydrate-cage volume changes, and composition-dependent thermal expansion.
A shared \(\theta_{\mc S}\) remains compatible with phase-specific thermal
expansion: each phase retains its own functions \(\mbf A_s^0\) and
\(\bar{\mbf F}_s^0\), and therefore its own derivative with respect to
\(\theta_{\mc S}\).
Thermo-poroelastic models make the same physical distinction by assigning pore
volume, fluid, and grain thermal expansions separately
\cite{palciauskas1982thermally,ghabezloo2009stress,ghabezloo2011micromechanics}.
As in standard finite-strain multiplicative plasticity
\cite{simo1988framework,simo1998} and Drumheller's mixture treatment
\cite{drumheller2000}, the true-plastic flow is assumed isochoric,
%
\begin{equation}
	\det\bar{\mbf F}_s^p=1,
	\qquad
	\operatorname{tr}\left[
		\dot{\bar{\mbf F}}_s^p
		(\bar{\mbf F}_s^p)^{-1}
		\right]=0,
	\qquad
	s\in\mc S .
	\label{eq:MC_solid_isochoric_plastic_flow}
\end{equation}
%
The dependent elastic arguments are
%
\begin{equation}
	\mbf A_s^e
	=
	\mbf A_s
	(\mbf A_s^0)^{-1}
	(\mbf A_s^p)^{-1},
	\qquad
	\bar{\mbf F}_s^e
	=
	\bar{\mbf F}_s
	(\bar{\mbf F}_s^0)^{-1}
	(\bar{\mbf F}_s^p)^{-1}.
	\label{eq:MC_solid_elastic_arguments}
\end{equation}
%

The solid component partial densities and mass fractions are
%
\begin{equation}
	\rho_s^\alpha
	=
	\rho_s\eta_s^\alpha
	=
	\phi_s\bar\rho_s\eta_s^\alpha,
	\qquad
	\sum_\alpha\eta_s^\alpha=1,
	\qquad
	\sum_\alpha\dot\eta_s^\alpha=0 .
	\label{eq:MC_solid_component_partial_density}
\end{equation}
%
When component intrinsic densities \(\bar\rho_s^\alpha\) and component solid
volume fractions \(\phi_s^\alpha\) are useful, they satisfy
%
\begin{subequations}
	\label{eq:MC_solid_density_mixture_rules}
	\begin{align}
		\bar\rho_s
		 & =
		\sum_\alpha
		\phi_s^\alpha\bar\rho_s^\alpha,
		\label{eq:MC_solid_intrinsic_density_rule} \\
		\eta_s^\alpha
		 & =
		\frac{\phi_s^\alpha\bar\rho_s^\alpha}{\bar\rho_s},
		\label{eq:MC_solid_mass_fraction_rule}
	\end{align}
	\begin{equation}
		\phi_s^\alpha
		=
		\frac{\eta_s^\alpha\bar\rho_s}{\bar\rho_s^\alpha},
		\qquad
		\sum_\alpha\phi_s^\alpha=1 .
		\label{eq:MC_solid_component_volume_fraction_rule}
	\end{equation}
\end{subequations}
%

The component material balances use the component material increment,
%
\begin{equation}
	J\phi_s\bar\rho_s\eta_s^\alpha
	=
	\phi_{s0}\bar\rho_{s0}\eta_{s0}^\alpha
	+\mc C_s^\alpha,
	\qquad
	s\in\mc S .
	\label{eq:MC_solid_component_material_mass}
\end{equation}
%
Summing over components gives
%
\begin{equation}
	J\phi_s\bar\rho_s
	=
	\phi_{s0}\bar\rho_{s0}+\mc C_s,
	\qquad
	\mc C_s=\sum_\alpha\mc C_s^\alpha .
	\label{eq:MC_solid_phase_material_mass}
\end{equation}
%
With true local conservation in the solid material,
%
\begin{equation}
	\bar J_s\bar\rho_s=\bar\rho_{s0},
	\label{eq:MC_solid_true_mass_conservation}
\end{equation}
%
\eqref{eq:MC_solid_phase_material_mass} becomes
%
\begin{equation}
	a_s\phi_s\bar\rho_{s0}
	=
	\phi_{s0}\bar\rho_{s0}+\mc C_s .
	\label{eq:MC_solid_distension_mass_relation}
\end{equation}
%
The rate form follows from \(J=a_s\bar J_s\),
%
\begin{empheq}[box=\fbox]{equation}
	\frac{\dot{a}_s}{a_s}
	=
	\nabla_{\mbf x}\cdot\mbf v_{\mc S}
	+\frac{\dot{\bar\rho}_s}{\bar\rho_s}.
	\label{eq:MC_solid_distension_evolution}
\end{empheq}
%
The intrinsic density rate needed in the Coleman--Noll analysis follows from
\eqref{eq:MC_solid_true_split}, isochoric plastic flow, and the
composition-dependent stress-free true deformation.  In the derivatives below,
the subscript \(\mc X_s\) indicates that all other members of the solid state
set defined in \eqref{eq:MC_state_sets} are held fixed:
%
\begin{align}
	\frac{\dot{\bar\rho}_s}{\bar\rho_s}
	={} &
	-\operatorname{tr}
	\left[
		\dot{\bar{\mbf F}}_s^e
		(\bar{\mbf F}_s^e)^{-1}
		\right]
	\notag  \\
	    & -
	\operatorname{tr}\left[
		\left(
		\left.
		\p{\bar{\mbf F}_s^0}{\theta_{\mc S}}
		\right\vert_{\mc X_s}
		\dot\theta_{\mc S}
		+
		\sum_\beta
		\left.
		\p{\bar{\mbf F}_s^0}{\eta_s^\beta}
		\right\vert_{\mc X_s}
		\dot\eta_s^\beta
		\right)
		(\bar{\mbf F}_s^0)^{-1}
		\right].
	\label{eq:MC_solid_density_rate}
\end{align}
%

The generic fluid and multicomponent solid Helmholtz energies are
%
\begin{equation}
	\begin{alignedat}{2}
		\psi_f
		& =
		\psi_f(
		\bar\rho_f,\{\eta_f^\alpha\},\theta_{\mc F}),
		\qquad&
		f\in\mc F,
		\\
		\psi_s
		& =
		\psi_s(
		\bar{\mbf F}_s^e,\mbf A_s^e,
		\bar\rho_s,\{\eta_s^\alpha\},\theta_{\mc S}),
		\qquad&
		s\in\mc S .
	\end{alignedat}
	\label{eq:MC_solid_general_free_energy}
\end{equation}
%
with concrete nonelectrical state sets
%
\begin{equation}
	\mc X_s
	=
	\{
	\bar{\mbf F}_s^e,\mbf A_s^e,
	\bar\rho_s,\{\eta_s^\alpha\},\theta_{\mc S}
	\},
	\qquad
	\mc X_f
	=
	\{
	\bar\rho_f,\{\eta_f^\alpha\},\theta_{\mc F}
	\}.
	\label{eq:MC_state_sets}
\end{equation}
%
The same held-fixed convention applies to derivatives annotated by
\(\mc X_f\).  The composition variables are constrained by their mass-fraction
normalizations, so composition
coefficients are projected onto the corresponding tangent spaces by subtracting
the \(N\)th component coefficient from the \(\beta=1,\ldots,N-1\) coefficients.

\pagebreak[3]
Electrical effects enter the thermodynamic structure through the phase
electric-enthalpy densities \(\omega^{+}_{\xi}\) introduced in
\sref{virtual_power_derivation}.  Each \(\omega_\xi^+\) is normalized per
current phase volume, and therefore contributes
\(\phi_\xi\omega_\xi^+\) per current mixture volume.  By comparison, the
interfacial Helmholtz free-energy density \(\gamma\) introduced in
\eqref{eq:interfacial_energy_density} is normalized per current fluid volume
and contributes \(\phi\gamma\) per current mixture volume.  The electric field
is a single spatial field, but each phase may have its own dielectric response
to that field.  The phase and total electric displacements are
%
\begin{equation}
	\mbf d_\xi
	=
	-\left.
	\p{\omega_\xi^+}{\mbf E}
	\right\vert_{\bar\rho_\xi,\{\eta_\xi^\alpha\},\theta_\xi},
	\qquad
	\mbf d
	=
	\sum_\xi\phi_\xi\mbf d_\xi .
	\label{eq:MC_electric_displacement_definitions}
\end{equation}
%
The electric field \(\mbf E\) is the independent electrical variable, whereas
the phase displacement \(\mbf d_\xi\) is the constitutive dielectric response.
This is the electric-enthalpy form of electroelastic potential theory
\cite{dorfmann2017nonlinear_electroelasticity}.

The fluid electrical enthalpy has the constitutive dependence admitted by the
variational state space,
\(\omega_f^+=\omega_f^+(\mbf E,\bar\rho_f,\{\eta_f^\alpha\},\theta_{\mc F})\).
Density, composition, and temperature therefore change the fluid permittivity
and its dielectric pressure.  Such state dependence is relevant across
reservoir lithologies: the bulk dielectric response of saturated sandstone
varies with porosity and pore-fluid salinity, while shale permittivity depends
additionally on water saturation, clay mineralogy, pore-fluid composition, and
bedding orientation
\cite{kirichek2018dielectric_sandstone,beloborodov2017dielectric_shales}.
The solid enthalpy
\(\omega_s^+\) has the analogous field, intrinsic-density, solid-composition,
and \(\theta_{\mc S}\) arguments.  These state arguments give the Maxwell stress
and the density, composition, and temperature contributions in the
corresponding Coleman--Noll restrictions.

The reservoir-scale constitutive choice takes \(\omega_\xi^+\) to be even in
\(\mbf E\).  Electrical coupling acts through Maxwell stress, phase equilibrium,
electrochemical transport, and electro-osmotic cross-coupling.  This focus is
consistent with saturated-rock self-potential measurements dominated by
electrokinetic response~\cite{moore2007self_potential} and tight-sandstone
corefloods in which electro-osmosis and electrochemical alteration govern the
electrical change in pressure demand, wettability, and oil
recovery~\cite{jia2024electrically_enhanced_oil_recovery}.  Plastic evolution
depends on the phase material stress and composition, so electrical loading
affects plastic response through the coupled Maxwell momentum solution.

The phase electrical stress \(\bs\sigma_\xi^{+}\) is the Maxwell stress defined in
\eqref{eq:vp_electrical_stress_definition}.  This stress is assigned to each
phase momentum coefficient and is included in the combined reversible stress
\(\bs\sigma_\xi^{\prime+}\).  The remaining density, composition, and thermal
derivatives of \(\omega^{+}_{\xi}\) are added to their matching reversible
coefficients before imposing the corresponding Coleman--Noll restrictions.

The generic Coleman--Noll argument uses the phase internal energy \(e_\xi\),
Helmholtz free energy \(\psi_\xi\), and entropy \(\mathfrak{s}_\xi\), with a
general composition-dependent phase energy.  The additive specialization in
\sref{MC_additive_component_specialization} resolves the solid Helmholtz energy
into component contributions.

\subsection{Coleman--Noll procedure}
\label{sec:MC_coleman_noll_procedure}

The Coleman--Noll procedure is the standard constitutive exploitation of the
local entropy inequality: the constitutive state dependence and balance laws are
substituted into the inequality, coefficients of independent admissible rates
are collected, and the remaining residual terms are required to be dissipative
\cite{coleman1974thermodynamics}.

The Hamilton--d'Alembert--Onsager construction in
\sref{virtual_power_derivation} and the Coleman--Noll procedure have
complementary roles.  The Hamilton--d'Alembert variations determine the phase
balance laws and constraint-multiplier relations, while the Onsager stationarity
condition supplies a near-equilibrium class of reaction and relative-transport
laws.  The Coleman--Noll procedure takes the balance laws and kinematics as given
and uses the entropy inequality to identify admissible constitutive responses and
force--rate relations.  Consequently, overlapping reversible relations provide
consistency checks when the same free energies and constraints are used.

Let \(e_\xi\) denote the specific internal energy of phase \(\xi\), and let
\(\mathfrak{s}_\xi\) denote its specific entropy.  For fluid phases,
%
\begin{equation}
	e_f=\psi_f+\theta_{\mc F}\mathfrak{s}_f,
	\qquad
	\psi_f=\psi_f(\bar\rho_f,\{\eta_f^\alpha\},\theta_{\mc F}).
	\label{eq:MC_fluid_legendre_relation}
\end{equation}
%
For multicomponent solid phases,
%
\begin{equation}
	e_s=\psi_s+\theta_{\mc S}\mathfrak{s}_s,
	\qquad
	\psi_s=\psi_s(
	\bar{\mbf F}_s^e,\mbf A_s^e,
	\bar\rho_s,\{\eta_s^\alpha\},\theta_{\mc S}).
	\label{eq:MC_solid_legendre_relation}
\end{equation}
%
The free-energy rates are
%
\begin{align}
	\sum_{f\in\mc F}\rho_f\dot\psi_f
	={} &
	\sum_{f\in\mc F}
	\rho_f
	\left(
	\left.
	\p{\psi_f}{\bar\rho_f}
	\right\vert_{\mc X_f}
	\dot{\bar\rho}_f
	+
	\sum_\beta
	\left.
	\p{\psi_f}{\eta_f^\beta}
	\right\vert_{\mc X_f}
	\dot\eta_f^\beta
	+
	\left.
		\p{\psi_f}{\theta_{\mc F}}
	\right\vert_{\mc X_f}
	\dot\theta_{\mc F}
	\right),
	\label{eq:MC_fluid_free_energy_rate}
	\\
	\sum_{s\in\mc S}\rho_s\dot\psi_s
	={} &
	\sum_{s\in\mc S}
	\rho_s
	\left[
		\left.
		\p{\psi_s}{\bar{\mbf F}_s^e}
		\right\vert_{\mc X_s}
		:
		\dot{\bar{\mbf F}}_s^e
		+
		\left.
		\p{\psi_s}{\mbf A_s^e}
		\right\vert_{\mc X_s}
		:
		\dot{\mbf A}_s^e
		\right.
	\notag       \\
	    & \left.
		+
		\left.
		\p{\psi_s}{\bar\rho_s}
		\right\vert_{\mc X_s}
		\dot{\bar\rho}_s
		+
		\sum_\beta
		\left.
		\p{\psi_s}{\eta_s^\beta}
		\right\vert_{\mc X_s}
		\dot\eta_s^\beta
		+
		\left.
		\p{\psi_s}{\theta_{\mc S}}
		\right\vert_{\mc X_s}
		\dot\theta_{\mc S}
		\right].
	\label{eq:MC_solid_free_energy_rate}
\end{align}
%
The fluid mechanical stress is
%
\begin{equation}
	\bs\sigma_f'=\mbf 0,
	\qquad
	f\in\mc F .
	\label{eq:MC_fluid_stress_free}
\end{equation}
%
The reversible fluid stress is therefore the Maxwell stress:
%
\begin{equation}
	\bs\sigma_f^{\prime+}
	=
	\bs{\sigma}^{+}_{f},
	\qquad
	f\in\mc F .
	\label{eq:MC_fluid_total_reversible_stress}
\end{equation}
%
The effective interfacial Helmholtz free-energy density \(\gamma\), normalized
per current fluid volume, retains the form introduced in
\eqref{eq:interfacial_energy_density}.  Its thermal argument is
\(\theta_{\mc F}\), because the saturations and capillary interfaces belong to
the mobile-fluid subsystem.  Allowing \(\gamma\) to depend on temperature is
consistent with standard capillary-pressure scaling used in reservoir
engineering \cite{leverett1941capillary}.  The collective density
\(\phi\gamma\) is stored in the skeleton-following representative volume.  Its
volume-fraction arguments therefore use
\eqref{eq:MC_skeleton_volume_fraction_rates}, and its thermal argument uses the
same skeleton rate of \(\theta_{\mc F}\).  Using the phase interfacial
potentials \(\gamma_f\) defined by
\eqref{eq:primitive_gamma_derivatives}, the skeleton-following chain rule gives
%
\begin{equation}
	\frac{D_{\mc S}}{Dt}
	\left[
		\phi\gamma
	\right]
	=
	\sum_{f\in\mc F}
	\gamma_f\dot\phi_f
	+
	\phi
	\left.
	\p{\gamma}{\theta_{\mc F}}
	\right\vert_{\{S_f\}_{f\in\mc F}}
	\dot\theta_{\mc F} .
	\label{eq:MC_interfacial_helmholtz_rate}
\end{equation}
%
Equation~\eqref{eq:MC_interfacial_helmholtz_rate} is a constitutive storage rate
in the summed mixture energy equation.  The present volume-averaged formulation
represents interfacial energy by the continuum storage density \(\phi\gamma\)
per current mixture volume and does not introduce an independently moving sharp
interface.  A sharp-interface extension would instead require an interface
velocity and separate surface balances of energy and entropy
\cite{bothe2022sharp_interface}.

Newly added component mass is assigned the specific internal energy \(e_\xi\)
of the phase receiving or losing that mass.
Let \(\mbf L_\xi=\nabla_{\mbf x}\mbf v_\xi\) be the phase
velocity gradient, \(\mbf q_\xi\) the heat flux, \(\rho_\xi r_\xi\) the external
heat supply per unit current mixture volume, and \(\rho_\xi e_\xi^{\rm int}\)
the internal energy supply from interphase exchange.

Following the constituent energy-balance structure of Drumheller and Bedford,
we postulate one energy balance for each thermal subsystem
\(\mc G\in\{\mc F,\mc S\}\)
\cite{drumheller1980thermomechanical}; its conversion-work terms use the
generalized form of Drumheller \cite[eq.~(50)]{drumheller2000}.  The fluid
balance contains the collective interfacial energy, and the solid balance
contains the solid-phase contributions.  Their sum is the total mixture energy
balance, so internal exchange between the two subsystems cancels.  All
thermodynamic storage rates in these balances use the skeleton derivative
\eqref{eq:MC_skeleton_rate_convention}.  This is the same control-volume choice
used by Seguin and Walkington, who write the energy and entropy balances on a
region convecting with the solid while retaining transport by fluid velocities
relative to the solid in the corresponding flux and interaction terms
\cite[eqs.~(10)--(12)]{seguin2019multicomponent}.  In the present formulation,
bulk relative motion remains in the phase momentum and interphase-power terms,
whereas component motion relative to a phase remains in the diffusion and
dispersion fluxes.

The electrical extension is written in
the electric-enthalpy representation used in thermo-electroelasticity
\cite{montanaro2015thermo_electroelasticity}, with Maxwell-stress accounting
consistent with electromagnetic energy-flux balances
\cite{nika2025heterogeneous_dielectrics}.  The electric enthalpy
\(\phi_\xi\omega_\xi^+\) per current mixture volume is attributed to phase
\(\xi\), whereas the interfacial Helmholtz contribution \(\phi\gamma\) is
shared and therefore enters the summed balance only once.  With the Maxwell
stress defined by \eqref{eq:vp_electrical_stress_definition}, the electrical
enthalpy is written in its skeleton-rate expansion so that its density,
composition, thermal, Maxwell-power, and electric-potential contributions are
visible term by term.  The postulated energy equation is
%
\begin{equation}
	\begin{aligned}
		 & \sum_{\xi\in\mc G}
		\left[
			\rho_\xi\dot e_\xi
			+
			\sum_\alpha\phasedot{c}_\xi^\alpha e_\xi
			+
			\omega_\xi^+\dot\phi_\xi
			+
			\phi_\xi
			\left.
			\p{\omega_\xi^+}{\bar\rho_\xi}
			\right\vert_{\mc X_\xi,\mbf E}
			\dot{\bar\rho}_\xi
			\right.
		\\
		 & \qquad
			\left.
			+
			\phi_\xi
			\sum_\alpha
			\left.
			\p{\omega_\xi^+}{\eta_\xi^\alpha}
			\right\vert_{\mc X_\xi,\mbf E}
			\dot\eta_\xi^\alpha
			+
			\phi_\xi
			\left.
			\p{\omega_\xi^+}{\theta_{\mc G}}
			\right\vert_{\mc X_\xi,\mbf E}
			\dot\theta_{\mc G}
			+
			\bs\sigma_\xi^+:\mbf L_\xi
			+
			\phi_\xi\mbf d_\xi\cdot\nabla_{\mbf x}\dot\varphi
			\right]
		\\
		 & \qquad
		+
		\delta_{\mc G\mc F}
		\left[
			\sum_{f\in\mc F}
			\gamma_f\dot\phi_f
			+
			\phi
			\left.
			\p{\gamma}{\theta_{\mc F}}
			\right\vert_{\{S_f\}_{f\in\mc F}}
			\dot\theta_{\mc F}
		\right]
		\\
		 & =
		\sum_{\xi\in\mc G}
		\left[
			\bs\sigma_\xi^{\prime+}:\mbf L_\xi
			+
			\phi_\xi\bar p_\xi
			\frac{\dot{\bar\rho}_\xi}{\bar\rho_\xi}
			+
			\bar p_\xi\dot\phi_\xi
			+
			\lambda\dot\phi_\xi
			+
			\sum_\alpha
			\frac{\pi_\xi^\alpha}{\eta_\xi^\alpha}
			\dot\eta_\xi^\alpha
			-
			\sum_\alpha L_\xi^\alpha\phasedot{c}_\xi^\alpha
			\right.
		\\
		 & \left.
			\qquad\qquad
			-
			\nabla_{\mbf x}\cdot\mbf q_\xi
			+
			\rho_\xi r_\xi
			+
			\rho_\xi e_\xi^{\rm int}
			\right].
	\end{aligned}
	\label{eq:MC_energy_balance}
\end{equation}
%
Here \(\delta_{\mc G\mc F}=1\) for the fluid balance and vanishes for the solid
balance, and
\(e_\xi=\psi_\xi+\theta_\xi\mathfrak{s}_\xi\).  The interfacial term is the
skeleton-following storage rate in
\eqref{eq:MC_interfacial_helmholtz_rate}.  The term
\(\lambda\sum_{\xi\in\mc G}\dot\phi_\xi\) is the reversible power exerted on
thermal subsystem \(\mc G\) by the pointwise volume constraint.  These powers
exchange energy between the two thermal subsystems and sum to zero in the total
mixture energy balance because \(\lambda\sum_\xi\dot\phi_\xi=0\) by
\eqref{eq:MC_skeleton_volume_fraction_rates}.  Placing this constraint power in
each subsystem energy balance before division by \(\theta_{\mc G}\) is essential
when \(\theta_{\mc F}\ne\theta_{\mc S}\).
The total stress
\(\bs\sigma_\xi^{\prime+}\) appears with positive mechanical power.  When the
explicit Maxwell power on the left of \eqref{eq:MC_energy_balance} is collected
with that total stress power, it leaves the nonelectrical material stress power
identified below.  The electric-potential-gradient term is integrated by parts
in the entropy exploitation; boundary electrical power is omitted.

The connection with a conventional thermo-electroelastic energy balance is most
transparent in the nonporous, single-phase limit.  Let
\(\mc G=\{\xi\}\), set \(\phi_\xi=1\), suppress conversion, composition,
interfacial, and interphase-exchange terms, and let the skeleton observer follow
phase \(\xi\).  Phase mass conservation then gives
\(\dot\rho_\xi/\rho_\xi=-\operatorname{tr}\mbf L_\xi\).  Using
\(\mbf E=-\nabla_{\mbf x}\varphi\),
\(\dot{\mbf E}=-\nabla_{\mbf x}\dot\varphi-\mbf L_\xi^T\mbf E\), and
\eqref{eq:vp_electrical_stress_definition}, the electrical terms in
\eqref{eq:MC_energy_balance} combine into the material-volume rate of
\(\omega_\xi^+\), and the balance reduces to
%
\begin{align}
	\rho_\xi
	\frac{D_\xi}{Dt}
	\left(
		e_\xi+\frac{\omega_\xi^+}{\rho_\xi}
	\right)
	={} &
	\rho_\xi\dot e_\xi
	+
	\dot\omega_\xi^+
	+
	\omega_\xi^+\operatorname{tr}\mbf L_\xi
	\notag \\
	={} &
	\left(
	\bs\sigma_\xi^{\prime+}
	-
	\bar p_\xi\mbf I
	\right):\mbf L_\xi
	-
	\nabla_{\mbf x}\cdot\mbf q_\xi
	+
	\rho_\xi r_\xi .
	\label{eq:MC_single_phase_energy_limit}
\end{align}
%
Reference~\cite{montanaro2015thermo_electroelasticity} instead writes its
equation~(12) for the material internal energy \(\varepsilon\), with electrical
power \(\mbf E^{M}\cdot\rho\dot{\bs\pi}\), where
\(\bs\pi=\mbf P/\rho\) is the polarization per unit mass.  Combining that
balance with its electric-enthalpy definition, equation~(16), gives
%
\begin{equation}
	\rho
	\frac{D}{Dt}
	\left(
	\varepsilon-\mbf E^{M}\cdot\bs\pi
	\right)
	=
	\bs\tau:\mbf L
	-
	\nabla_{\mbf x}\cdot\mbf q
	-
	\mbf P\cdot\dot{\mbf E}^{M}
	+
	\rho r .
	\label{eq:MC_montanaro_electric_enthalpy_limit}
\end{equation}
%
Thus \eqref{eq:MC_single_phase_energy_limit} and
\eqref{eq:MC_montanaro_electric_enthalpy_limit} use different allocations of
electrical energy and power.  The latter is a material balance with explicit
polarization work; the former includes the full displacement-based electric
potential and Maxwell stress.  Adding the electrostatic field-energy balance
and its associated boundary flux converts between these allocations.  No
independent polarization variable is required here: if an SI vacuum-field split
is desired for comparison, it is recovered from
\(\mbf P=\mbf d-\varepsilon_0\mbf E\).  Retaining \(\mbf d\) is more natural for
porous-media applications because Gauss' law is posed in terms of free charge,
while effective polarization caused by porosity, saturation, composition, and
temperature is already represented through the constitutive dependence of
\(\omega_\xi^+\).  A separate polarization evolution law would be needed only
for additional dielectric physics such as relaxation, hysteresis, or
ferroelectric response.

The solid stress power is decomposed by
\(\mbf F=\mbf A_s\bar{\mbf F}_s\),
%
\begin{equation}
	\mbf L_s
	=
	\dot{\mbf A}_s\mbf A_s^{-1}
	+
	\mbf A_s
	\dot{\bar{\mbf F}}_s
	\bar{\mbf F}_s^{-1}
	\mbf A_s^{-1}.
	\label{eq:MC_F_split_for_power}
\end{equation}
%
Using \eqref{eq:MC_solid_kinematic_splits}, the solid stress power is
%
\begin{align}
	\sum_{s\in\mc S}\bs\sigma_s^{\prime+}:\mbf L_s
	={} &
	\sum_{s\in\mc S}
	\operatorname{tr}\left[
		(\bs\sigma_s^{\prime+})^T
		\dot{\mbf A}_s^e
		(\mbf A_s^e)^{-1}
		\right]
	\notag  \\
	    & +
	\sum_{s\in\mc S}
	\operatorname{tr}\left[
		(\mbf A_s^e)^{-1}
		(\bs\sigma_s^{\prime+})^T
		\mbf A_s^e
		\dot{\mbf A}_s^p
		(\mbf A_s^p)^{-1}
		\right]
	\notag  \\
	    & +
	\sum_{s\in\mc S}
	\operatorname{tr}\left[
		(\mbf A_s^p)^{-1}
		(\mbf A_s^e)^{-1}
		(\bs\sigma_s^{\prime+})^T
		\mbf A_s^e\mbf A_s^p
		\left.
		\p{\mbf A_s^0}{\theta_{\mc S}}
		\right\vert_{\mc X_s}
		(\mbf A_s^0)^{-1}
		\right]\dot\theta_{\mc S}
	\notag  \\
	    & +
	\sum_{s\in\mc S}
	\sum_\beta
	\operatorname{tr}\left[
		(\mbf A_s^p)^{-1}
		(\mbf A_s^e)^{-1}
		(\bs\sigma_s^{\prime+})^T
		\mbf A_s^e\mbf A_s^p
		\left.
		\p{\mbf A_s^0}{\eta_s^\beta}
		\right\vert_{\mc X_s}
		(\mbf A_s^0)^{-1}
		\right]\dot\eta_s^\beta
	\notag  \\
	    & +
	\sum_{s\in\mc S}
	\operatorname{tr}\left[
	\mbf A_s^{-1}
	(\bs\sigma_s^{\prime+})^T
	\mbf A_s
	\dot{\bar{\mbf F}}_s^e
	(\bar{\mbf F}_s^e)^{-1}
	\right]
	\notag  \\
	    & +
	\sum_{s\in\mc S}
	\operatorname{tr}\left[
	(\bar{\mbf F}_s^e)^{-1}
	\mbf A_s^{-1}
	(\bs\sigma_s^{\prime+})^T
	\mbf A_s
	\bar{\mbf F}_s^e
	\dot{\bar{\mbf F}}_s^p
	(\bar{\mbf F}_s^p)^{-1}
	\right]
	\notag  \\
	    & +
	\sum_{s\in\mc S}
	\operatorname{tr}\left[
	(\bar{\mbf F}_s^p)^{-1}
	(\bar{\mbf F}_s^e)^{-1}
	\mbf A_s^{-1}
	(\bs\sigma_s^{\prime+})^T
	\mbf A_s
	\bar{\mbf F}_s^e
	\bar{\mbf F}_s^p
	\left.
	\p{\bar{\mbf F}_s^0}{\theta_{\mc S}}
	\right\vert_{\mc X_s}
	(\bar{\mbf F}_s^0)^{-1}
	\right]\dot\theta_{\mc S}
	\notag  \\
	    & +
	\sum_{s\in\mc S}
	\sum_\beta
	\operatorname{tr}\left[
	(\bar{\mbf F}_s^p)^{-1}
	(\bar{\mbf F}_s^e)^{-1}
	\mbf A_s^{-1}
	(\bs\sigma_s^{\prime+})^T
	\mbf A_s
	\bar{\mbf F}_s^e
	\bar{\mbf F}_s^p
	\left.
	\p{\bar{\mbf F}_s^0}{\eta_s^\beta}
	\right\vert_{\mc X_s}
	(\bar{\mbf F}_s^0)^{-1}
	\right]\dot\eta_s^\beta .
	\label{eq:MC_stress_power_split}
\end{align}
%

Following the standard Liu multiplier exploitation of constrained entropy
inequalities \cite{liu1972method}, the entropy inequality is written in the same
skeleton observer frame,
%
\begin{align}
	0\le{} &
	\sum_{\mc G\in\{\mc F,\mc S\}}
	\sum_{\xi\in\mc G}
	\left[
		\rho_\xi\dot{\mathfrak{s}}_\xi
		+
		\sum_\alpha
		\phasedot{c}_\xi^\alpha
		\mathfrak{s}_\xi
		+
		\nabla_{\mbf x}\cdot
		\left(
		\frac{\mbf q_\xi}{\theta_{\mc G}}
		\right)
		-
		\frac{\rho_\xi r_\xi}{\theta_{\mc G}}
		\right].
	\label{eq:MC_entropy_inequality}
\end{align}
%
Eliminating the external heat supply in each thermal group with
\eqref{eq:MC_energy_balance} and adding the identically vanishing
mass-fraction tangent terms
\(\Pi_\xi\sum_\alpha\dot\eta_\xi^\alpha\) gives the group-resolved inequality
below.  The conservative electric-field identity and Gauss' law are imposed
before this thermal collection, leaving the material stress
\(\bs\sigma_\xi'\) and the density, composition, and temperature derivatives
of \(\omega_\xi^+\).  This avoids assigning the single electrostatic field a
separate subsystem temperature.
%
\begin{align}
	0\le{} &
	\sum_{\mc G\in\{\mc F,\mc S\}}
	\sum_{\xi\in\mc G}
	\begin{aligned}[t]
		 & \frac{1}{\theta_{\mc G}}
		\left(
		\bs\sigma_\xi':\mbf L_\xi
		+
		\phi_\xi\bar p_\xi
		\frac{\dot{\bar\rho}_\xi}{\bar\rho_\xi}
		-
		\rho_\xi
		\left(
		\dot\psi_\xi+\mathfrak{s}_\xi\dot\theta_{\mc G}
		\right)
		-
		\sum_\alpha
		\phasedot{c}_\xi^\alpha
		\left(
		\psi_\xi+L_\xi^\alpha
		\right)
		+
		\rho_\xi e_\xi^{\rm int}
		\right)
	\end{aligned}
	\notag          \\
	       & -
	\sum_{\mc G\in\{\mc F,\mc S\}}
	\sum_{\xi\in\mc G}
	\frac{
	\mbf q_\xi\cdot\nabla_{\mbf x}\theta_{\mc G}
	}{
	\theta_{\mc G}^2
	}
	\notag          \\
	       & +
	\sum_{\mc G\in\{\mc F,\mc S\}}
		\frac{1}{\theta_{\mc G}}
		\sum_{\xi\in\mc G}
		\left(
		\left(
		\lambda+\bar p_\xi
		\right)\dot\phi_\xi
		+
		\sum_\alpha
	\frac{\pi_\xi^\alpha}{\eta_\xi^\alpha}
	\dot\eta_\xi^\alpha
	+
	\Pi_\xi\sum_\alpha\dot\eta_\xi^\alpha
	\right)
	\notag          \\
	       & -
	\frac{1}{\theta_{\mc F}}
	\left[
		\sum_{f\in\mc F}
		\gamma_f\dot\phi_f
		+
		\phi
		\left.
		\p{\gamma}{\theta_{\mc F}}
		\right\vert_{\{S_f\}_{f\in\mc F}}
		\dot\theta_{\mc F}
	\right]
	\notag          \\
	       & +
	\sum_{\mc G\in\{\mc F,\mc S\}}
	\frac{1}{\theta_{\mc G}}
	\left[
		-
		\sum_{\xi\in\mc G}
		\left(
		\omega^{+}_{\xi}\dot\phi_\xi
		+
		\phi_\xi
		\left.
		\p{\omega_\xi^+}{\bar\rho_\xi}
		\right\vert_{\mc X_\xi,\mbf E}
		\dot{\bar\rho}_\xi
		\right)
		\right.
	\notag          \\
	       & \left.
		-
		\sum_{\xi\in\mc G}
		\phi_\xi
		\sum_\beta
		\left.
		\p{\omega_\xi^+}{\eta_\xi^\beta}
		\right\vert_{\mc X_\xi,\mbf E}
		\dot\eta_\xi^\beta
		-
		\sum_{\xi\in\mc G}
		\phi_\xi
		\left.
		\p{\omega_\xi^+}{\theta_{\mc G}}
		\right\vert_{\mc X_\xi,\mbf E}
		\dot\theta_{\mc G}
		\right].
	\label{eq:MC_reduced_entropy_inequality}
\end{align}
%

The interfacial term in \eqref{eq:MC_reduced_entropy_inequality} is the negative
of the skeleton-following storage rate in
\eqref{eq:MC_interfacial_helmholtz_rate}.  The electric-enthalpy derivatives
come directly from \eqref{eq:MC_energy_balance}.  The conservative field-rate
and Maxwell-power terms have already been removed with Gauss' law and
\begin{equation}
	\bs\sigma_\xi^{\prime+}:\mbf L_\xi
	-
	\bs\sigma_\xi^+:\mbf L_\xi
	=
	\bs\sigma_\xi':\mbf L_\xi .
	\label{eq:MC_material_stress_power_identity}
\end{equation}
This identity identifies \(\bs\sigma_\xi'\) as the stress power conjugate to
the phase Helmholtz-energy rate.  The momentum equation retains the total
stress \(\bs\sigma_\xi^{\prime+}\).  Poynting-flux derivations give the same
allocation: the first law contains total stress power and field-energy transport
supplies the Maxwell-stress power in the Clausius--Duhem inequality
\cite{nika2025heterogeneous_dielectrics}.  Equivalently, an
electric-enthalpy Legendre transform gives the displacement with a minus
potential derivative and retains the total stress as the deformation
derivative of the combined potential
\cite{montanaro2015thermo_electroelasticity,
	dorfmann2017nonlinear_electroelasticity}.

Substituting \eqref{eq:MC_fluid_free_energy_rate}, the explicit
electrical-enthalpy terms,
\eqref{eq:MC_solid_free_energy_rate}, \eqref{eq:MC_solid_density_rate}, and
\eqref{eq:MC_stress_power_split}, while using
\eqref{eq:MC_material_stress_power_identity}, gives the collected inequality,
%
	\begin{align}
		0\le{} &
		\frac{1}{\theta_{\mc S}}
		\sum_{s\in\mc S}
	\operatorname{tr}\left[
		\left(
		(\bs\sigma_s')^T
		-
		\rho_s
		\mbf A_s^e
		\left(
		\left.
		\p{\psi_s}{\mbf A_s^e}
		\right\vert_{\mc X_s}
		\right)^T
		\right)
		\dot{\mbf A}_s^e
		(\mbf A_s^e)^{-1}
		\right]
	\notag                  \\
		       & +
		\frac{1}{\theta_{\mc S}}
		\sum_{s\in\mc S}
	\operatorname{tr}\left[
	\left(
	\mbf A_s^{-1}
	(\bs\sigma_s')^T
	\mbf A_s
	-
	\rho_s
	\bar{\mbf F}_s^e
	\left(
	\left.
	\p{\psi_s}{\bar{\mbf F}_s^e}
	\right\vert_{\mc X_s}
	\right)^T
	\right.
	\right.
	\notag                  \\
	       & \left.
	\left.
	+
	\left(
	\rho_s
	\left.
	\p{\psi_s}{\bar\rho_s}
	\right\vert_{\mc X_s}
	\bar\rho_s
	+
	\phi_s
	\left.
	\p{\omega^{+}_{s}}{\bar\rho_s}
	\right\vert_{\mc X_s,\mbf E}
	\bar\rho_s
	-
	\phi_s\bar p_s
	\right)\mbf I
	\right)
	\dot{\bar{\mbf F}}_s^e
	(\bar{\mbf F}_s^e)^{-1}
	\right]
	\notag                  \\
		       & +
		\frac{1}{\theta_{\mc S}}
		\sum_{s\in\mc S}
	\operatorname{tr}\left[
		(\mbf A_s^e)^{-1}
		(\bs\sigma_s')^T
		\mbf A_s^e
		\dot{\mbf A}_s^p
		(\mbf A_s^p)^{-1}
		\right]
	\notag                  \\
		       & +
		\frac{1}{\theta_{\mc S}}
		\sum_{s\in\mc S}
	\operatorname{tr}\left[
	(\bar{\mbf F}_s^e)^{-1}
	\mbf A_s^{-1}
	(\bs\sigma_s')^T
	\mbf A_s
	\bar{\mbf F}_s^e
	\dot{\bar{\mbf F}}_s^p
	(\bar{\mbf F}_s^p)^{-1}
	\right]
	\notag                  \\
		       & +
		\frac{1}{\theta_{\mc F}}
		\sum_{f\in\mc F}
	\left[
		\frac{\phi_f\bar p_f}{\bar\rho_f}
		-
		\rho_f
		\left.
		\p{\psi_f}{\bar\rho_f}
		\right\vert_{\mc X_f}
		-
		\phi_f
		\left.
		\p{\omega^{+}_{f}}{\bar\rho_f}
		\right\vert_{\mc X_f,\mbf E}
		\right]
	\dot{\bar\rho}_f
	\notag                  \\
		       & +
		\frac{1}{\theta_{\mc F}}
		\sum_{f\in\mc F}
	\sum_\beta
	\left[
		\frac{\pi_f^\beta}{\eta_f^\beta}
		+\Pi_f
		-
		\rho_f
		\left.
		\p{\psi_f}{\eta_f^\beta}
		\right\vert_{\mc X_f}
		-
		\phi_f
		\left.
		\p{\omega^{+}_{f}}{\eta_f^\beta}
		\right\vert_{\mc X_f,\mbf E}
		\right]\dot\eta_f^\beta
	\notag \displaybreak[4] \\
		       & +
		\frac{1}{\theta_{\mc S}}
		\sum_{s\in\mc S}
	\sum_\beta
	\left[
	\frac{\pi_s^\beta}{\eta_s^\beta}
	+\Pi_s
	-
	\rho_s
	\left.
	\p{\psi_s}{\eta_s^\beta}
	\right\vert_{\mc X_s}
	-
	\phi_s
	\left.
	\p{\omega^{+}_{s}}{\eta_s^\beta}
	\right\vert_{\mc X_s,\mbf E}
	\right.
	\notag                  \\
	       & \left.
	+
	\operatorname{tr}\left[
		(\mbf A_s^p)^{-1}
		(\mbf A_s^e)^{-1}
		(\bs\sigma_s')^T
		\mbf A_s^e\mbf A_s^p
		\left.
		\p{\mbf A_s^0}{\eta_s^\beta}
		\right\vert_{\mc X_s}
		(\mbf A_s^0)^{-1}
		\right]
	\right.
	\notag                  \\
	       & \left.
	+
	\operatorname{tr}\left[
	\left(
	(\bar{\mbf F}_s^p)^{-1}
	(\bar{\mbf F}_s^e)^{-1}
	\mbf A_s^{-1}
	(\bs\sigma_s')^T
	\mbf A_s
	\bar{\mbf F}_s^e
	\bar{\mbf F}_s^p
	\right.
	\right.
	\right.
	\notag                  \\
	       & \left.
	\left.
	\left.
	+
	\left(
	\rho_s
	\left.
	\p{\psi_s}{\bar\rho_s}
	\right\vert_{\mc X_s}
	\bar\rho_s
	+
	\phi_s
	\left.
	\p{\omega^{+}_{s}}{\bar\rho_s}
	\right\vert_{\mc X_s,\mbf E}
	\bar\rho_s
	-
	\phi_s\bar p_s
	\right)\mbf I
	\right)
	\left.
	\p{\bar{\mbf F}_s^0}{\eta_s^\beta}
	\right\vert_{\mc X_s}
	(\bar{\mbf F}_s^0)^{-1}
	\right]
	\right]\dot\eta_s^\beta
	\notag                  \\
		       & +
		\frac{1}{\theta_{\mc F}}
		\left[
		-\sum_{f\in\mc F}
	\rho_f
	\left(
	\mathfrak{s}_f
	+
	\left.
		\p{\psi_f}{\theta_{\mc F}}
	\right\vert_{\mc X_f}
	\right)
	-
	\sum_{f\in\mc F}
	\phi_f
	\left.
		\p{\omega^{+}_{f}}{\theta_{\mc F}}
		\right\vert_{\mc X_f,\mbf E}
		-
		\phi
		\left.
		\p{\gamma}{\theta_{\mc F}}
		\right\vert_{\{S_f\}_{f\in\mc F}}
		\right]\dot\theta_{\mc F}
		\notag                  \\
		       & +
		\frac{1}{\theta_{\mc S}}
		\left[
		-
		\sum_{s\in\mc S}
	\rho_s
	\left(
	\mathfrak{s}_s
	+
	\left.
		\p{\psi_s}{\theta_{\mc S}}
	\right\vert_{\mc X_s}
	\right)
	-
	\sum_{s\in\mc S}
	\phi_s
	\left.
		\p{\omega^{+}_{s}}{\theta_{\mc S}}
		\right\vert_{\mc X_s,\mbf E}
		\right.
	\notag                  \\
	       & \left.
	+
	\sum_{s\in\mc S}
	\operatorname{tr}\left[
		(\mbf A_s^p)^{-1}
		(\mbf A_s^e)^{-1}
		(\bs\sigma_s')^T
		\mbf A_s^e\mbf A_s^p
		\left.
		\p{\mbf A_s^0}{\theta_{\mc S}}
		\right\vert_{\mc X_s}
		(\mbf A_s^0)^{-1}
		\right]
	\right.
	\notag                  \\
	       & \left.
	+
	\sum_{s\in\mc S}
	\operatorname{tr}\left[
	\left(
	(\bar{\mbf F}_s^p)^{-1}
	(\bar{\mbf F}_s^e)^{-1}
	\mbf A_s^{-1}
	(\bs\sigma_s')^T
	\mbf A_s
	\bar{\mbf F}_s^e
	\bar{\mbf F}_s^p
	\right.
	\right.
	\right.
	\notag                  \\
	       & \left.
	\left.
	\left.
	+
	\left(
	\rho_s
	\left.
	\p{\psi_s}{\bar\rho_s}
	\right\vert_{\mc X_s}
	\bar\rho_s
	+
	\phi_s
	\left.
	\p{\omega^{+}_{s}}{\bar\rho_s}
	\right\vert_{\mc X_s,\mbf E}
	\bar\rho_s
	-
	\phi_s\bar p_s
	\right)\mbf I
	\right)
	\left.
		\p{\bar{\mbf F}_s^0}{\theta_{\mc S}}
	\right\vert_{\mc X_s}
	(\bar{\mbf F}_s^0)^{-1}
	\right]
		\right]\dot\theta_{\mc S}
	\notag                  \\
	       & +
	\frac{1}{\theta_{\mc S}}
		\sum_{s\in\mc S}
		\left(
		\lambda
		+
		\bar p_s
		-
		\omega^{+}_{s}
		\right)\dot\phi_s
	+
	\frac{1}{\theta_{\mc F}}
		\sum_{f\in\mc F}
		\left(
		\lambda
		+
		\bar p_f
	-
	\omega^{+}_{f}
	-
	\gamma_f
	\right)\dot\phi_f
	\notag                  \\
	       & -
	\sum_m
	\left[
	\frac{1}{\theta_{\mc F}}
	\sum_{f\in\mc F}
	\sum_\alpha
	\nu_{f (m)}^\alpha
	\left(
		\psi_f+L_f^\alpha
	\right)
	+
	\frac{1}{\theta_{\mc S}}
	\sum_{s\in\mc S}
	\sum_\alpha
	\nu_{s (m)}^\alpha
	\left(
		\psi_s+L_s^\alpha
		\right)
	\right]\dot r_{(m)}
	\notag                  \\
	       & -
	\sum_{\mc G\in\{\mc F,\mc S\}}
	\sum_{\xi\in\mc G}\sum_\alpha
	\left(
	\mbf j_{\xi,{\rm disp}}^\alpha
	+
	\mbf j_{\xi,{\rm diff}}^\alpha
	\right)\cdot
	\nabla_{\mbf x}\left(
	\frac{\psi_\xi+L_\xi^\alpha}{\theta_{\mc G}}
	\right)
	\notag                  \\
	       & +
	\sum_{\mc G\in\{\mc F,\mc S\}}
	\frac{1}{\theta_{\mc G}}
	\sum_{\xi\in\mc G}
	\rho_\xi e_\xi^{\rm int}
	-
	\sum_{\mc G\in\{\mc F,\mc S\}}
	\sum_{\xi\in\mc G}
	\frac{
	\mbf q_\xi\cdot\nabla_{\mbf x}\theta_{\mc G}
	}{
	\theta_{\mc G}^2
	}
	.
	\label{eq:MC_collected_entropy}
\end{align}
%

\subsection{Reversible restrictions}
\label{sec:MC_reversible_restrictions}

The Coleman--Noll collection uses the material stress
\(\bs\sigma_\xi'\) identified by
\eqref{eq:MC_material_stress_power_identity}.  The total stress
\(\bs\sigma_\xi^{\prime+}\) is the sum of the phase Helmholtz stress
\(\bs\sigma_\xi'\) and the Maxwell stress \(\bs\sigma_\xi^+\).

The coefficient of
\(\dot{\bar{\mbf F}}_s^e(\bar{\mbf F}_s^e)^{-1}\) gives
%
\begin{align}
	\mbf A_s^{-1}
	(\bs\sigma_s')^T
	\mbf A_s
	={} &
	\rho_s
	\bar{\mbf F}_s^e
	\left(
	\left.
	\p{\psi_s}{\bar{\mbf F}_s^e}
	\right\vert_{\mc X_s}
	\right)^T
	\notag  \\
	    & -
	\left(
	\rho_s
	\left.
	\p{\psi_s}{\bar\rho_s}
	\right\vert_{\mc X_s}
	\bar\rho_s
	+
	\phi_s
	\left.
	\p{\omega^{+}_{s}}{\bar\rho_s}
	\right\vert_{\mc X_s,\mbf E}
	\bar\rho_s
	-
	\phi_s\bar p_s
	\right)\mbf I .
	\label{eq:MC_stress_before_pressure}
\end{align}
%
For a fluid, \(\bs\sigma_f'=\mbf0\) by
\eqref{eq:MC_fluid_stress_free}, so
\(\bs\sigma_f^{\prime+}=\bs\sigma_f^+\).  Its reversible power is paired with
the electric-enthalpy rate in the entropy inequality, and the stress enters the
fluid momentum balance \eqref{eq:el_mom_f_simplified}.

The density coefficients identify the phase pressures,
%
\begin{subequations}
	\label{eq:MC_phase_pressure_definitions}
	\begin{align}
		\bar p_f
		 & =
		\bar\rho_f^2
		\left.
		\p{\psi_f}{\bar\rho_f}
		\right\vert_{\mc X_f}
		+
		\bar\rho_f
		\left.
		\p{\omega^{+}_{f}}{\bar\rho_f}
		\right\vert_{\mc X_f,\mbf E},
		\qquad
		f\in\mc F,
		\label{eq:MC_fluid_pressure_definition} \\
		\bar p_s
		 & =
		\bar\rho_s^2
		\left.
		\p{\psi_s}{\bar\rho_s}
		\right\vert_{\mc X_s}
		+
		\bar\rho_s
		\left.
		\p{\omega^{+}_{s}}{\bar\rho_s}
		\right\vert_{\mc X_s,\mbf E},
		\qquad
		s\in\mc S .
		\label{eq:MC_solid_pressure_definition}
	\end{align}
\end{subequations}
%
Substitution of \eqref{eq:MC_solid_pressure_definition} into
\eqref{eq:MC_stress_before_pressure} gives the material solid stress,
%
\begin{equation}
	\bs\sigma_s'
	=
	\mbf A_s^{-T}
	\left(
	\rho_s
	\left.
	\p{\psi_s}{\bar{\mbf F}_s^e}
	\right\vert_{\mc X_s}
	(\bar{\mbf F}_s^e)^T
	\right)
	\mbf A_s^T .
	\label{eq:MC_solid_stress_definition}
\end{equation}
%
The reversible electrical contribution derived in
\eqref{eq:vp_electrical_mechanical_power} assigns the complete electrical stress
to the solid phase when the solid phase map is varied:
%
\begin{equation}
	\bs{\sigma}^{+}_{s}
	=
	\left.
	\bs{\sigma}^{+}_{\xi}
	\right\vert_{\xi=s},
	\label{eq:MC_solid_electrical_stress_contribution}
\end{equation}
%
Equation~\eqref{eq:vp_electrical_stress_definition} identifies
\(\bs\sigma_s^+\) as the phase-attributed Maxwell stress for the adopted
electric-enthalpy state variables.

The total solid stress is recovered from
\eqref{eq:MC_material_stress_power_identity} as the sum of the material stress
and the complete electrical stress.

For the linear dielectric enthalpy
\(\omega^{+}_{\xi}=-\frac{1}{2}\mbf E\cdot\mbf d_\xi\), the sum of the
phase-attributed electrical stresses is the usual Maxwell stress
\(\mbf E\otimes\mbf d-\frac{1}{2}\left(\mbf E\cdot\mbf d\right)\mbf I\), with
\(\mbf d=\sum_\xi\phi_\xi\mbf d_\xi\).

The reversible distension coefficient gives
%
\begin{equation}
	\bs\sigma_s'
	=
	\rho_s
	\left.
	\p{\psi_s}{\mbf A_s^e}
	\right\vert_{\mc X_s}
	(\mbf A_s^e)^T .
	\label{eq:MC_distension_force_definition}
\end{equation}
%
The volume-fraction coefficients in
\eqref{eq:MC_collected_entropy} reproduce the reversible variational equations.
The constraint power included in each subsystem energy balance places the same
multiplier \(\lambda\) inside both temperature-weighted coefficients.  Because
each temperature is positive, the coefficients vanish precisely when
%
\begin{subequations}
	\label{eq:MC_volume_fraction_restrictions}
	\begin{align}
		0
		      & =
		\lambda
		+
		\bar p_s
		-
		\omega^{+}_{s},
		\quad &
		s\in\mc S,
		\label{eq:MC_solid_volume_fraction_restriction} \\
		0
		      & =
		\lambda
		+
		\bar p_f
		-
		\omega^{+}_{f}
		-\gamma_f,
		\quad &
		f\in\mc F .
		\label{eq:MC_fluid_volume_fraction_restriction}
	\end{align}
\end{subequations}
%
Thus the subsystem constraint powers give the same cancellation for
\(\theta_{\mc F}\ne\theta_{\mc S}\).

The fluid composition projection is
%
\begin{empheq}[box=\fbox]{equation}
	\rho_f
	\left(
	\left.
	\p{\psi_f}{\eta_f^\beta}
	\right\vert_{\mc X_f}
	-
	\left.
	\p{\psi_f}{\eta_f^N}
	\right\vert_{\mc X_f}
	\right)
	+
	\phi_f
	\left(
	\left.
	\p{\omega^{+}_{f}}{\eta_f^\beta}
	\right\vert_{\mc X_f,\mbf E}
	-
	\left.
	\p{\omega^{+}_{f}}{\eta_f^N}
	\right\vert_{\mc X_f,\mbf E}
	\right)
	=
	\frac{\pi_f^\beta}{\eta_f^\beta}
	-
	\frac{\pi_f^N}{\eta_f^N},
	\qquad
	\beta=1,\ldots,N-1.
	\label{eq:MC_fluid_composition_projection}
\end{empheq}
%
The multicomponent solid composition projection is
%
\begin{empheq}[box=\fbox]{align}
	& \rho_s
	\left(
	\left.
	\p{\psi_s}{\eta_s^\beta}
	\right\vert_{\mc X_s}
	-
	\left.
	\p{\psi_s}{\eta_s^N}
	\right\vert_{\mc X_s}
	\right)
	+
	\phi_s
	\left(
	\left.
	\p{\omega^{+}_{s}}{\eta_s^\beta}
	\right\vert_{\mc X_s,\mbf E}
	-
	\left.
	\p{\omega^{+}_{s}}{\eta_s^N}
	\right\vert_{\mc X_s,\mbf E}
	\right)
	\notag    \\
	& \quad
	-
	\operatorname{tr}\left[
		(\mbf A_s^p)^{-1}
		(\mbf A_s^e)^{-1}
		(\bs\sigma_s')^T
		\mbf A_s^e\mbf A_s^p
		\left(
		\left.
		\p{\mbf A_s^0}{\eta_s^\beta}
		\right\vert_{\mc X_s}
		-
		\left.
		\p{\mbf A_s^0}{\eta_s^N}
		\right\vert_{\mc X_s}
		\right)
		(\mbf A_s^0)^{-1}
		\right]
	\notag    \\
	& \quad
	-
	\operatorname{tr}\left[
	\left(
	(\bar{\mbf F}_s^p)^{-1}
	(\bar{\mbf F}_s^e)^{-1}
	\mbf A_s^{-1}
	(\bs\sigma_s')^T
	\mbf A_s
	\bar{\mbf F}_s^e
	\bar{\mbf F}_s^p
	\right)
	\right.
	\notag    \\
	& \left.
	\quad
	\left(
	\left.
	\p{\bar{\mbf F}_s^0}{\eta_s^\beta}
	\right\vert_{\mc X_s}
	-
	\left.
	\p{\bar{\mbf F}_s^0}{\eta_s^N}
	\right\vert_{\mc X_s}
	\right)
	(\bar{\mbf F}_s^0)^{-1}
	\right]
	\notag    \\
	& \qquad
	=
	\frac{\pi_s^\beta}{\eta_s^\beta}
	-
	\frac{\pi_s^N}{\eta_s^N},
	\qquad
	\beta=1,\ldots,N-1.
	\label{eq:MC_solid_composition_projection}
\end{empheq}
%
The two independent temperature rates give separate fluid and solid entropy
restrictions.  The interfacial contribution belongs to the fluid subsystem
because \(\gamma\) depends on \(\theta_{\mc F}\):
%
\begin{subequations}
	\label{eq:MC_entropy_definitions}
	\begin{align}
		\sum_{f\in\mc F}\rho_f\mathfrak{s}_f
		={} &
		-
		\sum_{f\in\mc F}\rho_f
		\left.
		\p{\psi_f}{\theta_{\mc F}}
		\right\vert_{\mc X_f}
		-
		\sum_{f\in\mc F}\phi_f
		\left.
		\p{\omega^{+}_{f}}{\theta_{\mc F}}
		\right\vert_{\mc X_f,\mbf E}
		-
		\phi
		\left.
		\p{\gamma}{\theta_{\mc F}}
		\right\vert_{\{S_f\}_{f\in\mc F}},
		\label{eq:MC_fluid_entropy_definition}
		\\
		\sum_{s\in\mc S}\rho_s\mathfrak{s}_s
		={} &
		-
		\sum_{s\in\mc S}\rho_s
		\left.
		\p{\psi_s}{\theta_{\mc S}}
		\right\vert_{\mc X_s}
		-
		\sum_{s\in\mc S}\phi_s
		\left.
		\p{\omega^{+}_{s}}{\theta_{\mc S}}
		\right\vert_{\mc X_s,\mbf E}
		\notag  \\
		    & +
		\sum_{s\in\mc S}
		\operatorname{tr}\left[
			(\mbf A_s^p)^{-1}
			(\mbf A_s^e)^{-1}
			(\bs\sigma_s')^T
			\mbf A_s^e\mbf A_s^p
			\left.
			\p{\mbf A_s^0}{\theta_{\mc S}}
			\right\vert_{\mc X_s}
			(\mbf A_s^0)^{-1}
			\right]
		\notag  \\
		    & +
		\sum_{s\in\mc S}
		\operatorname{tr}\left[
		(\bar{\mbf F}_s^p)^{-1}
		(\bar{\mbf F}_s^e)^{-1}
		\mbf A_s^{-1}
		(\bs\sigma_s')^T
		\mbf A_s
		\bar{\mbf F}_s^e
		\bar{\mbf F}_s^p
		\left.
		\p{\bar{\mbf F}_s^0}{\theta_{\mc S}}
		\right\vert_{\mc X_s}
		(\bar{\mbf F}_s^0)^{-1}
		\right].
		\label{eq:MC_solid_entropy_definition}
	\end{align}
\end{subequations}
%

\subsection{Sources, residual dissipation, and closures}
\label{sec:MC_residual_dissipation}

\subsubsection{Component supply and thermodynamic forces}
\label{sec:MC_component_supply_forces}

Converted and relatively transported mass of the same component enter the phase
material measure through their sum.  The variational derivation leaves the
generalized work potential \(L_\xi^\alpha\) explicit in
\eqref{eq:el_conversion_component}.  The entropy inequality identifies the
electrochemical potential as the thermodynamic conjugate to component supply.
Accordingly, set
%
\begin{equation}
	L_\xi^\alpha
	=
	\mu_\xi^\alpha-\psi_\xi .
	\label{eq:MC_L_mu_psi_closure}
\end{equation}
%
Then \(\psi_\xi+L_\xi^\alpha=\mu_\xi^\alpha\), with
\(\mu_\xi^\alpha\) interpreted as the electrochemical potential when charge is
retained.
Substituting \eqref{eq:MC_L_mu_psi_closure} into the coefficient of
\(\delta\mc C_\xi^\alpha\) in \eqref{eq:el_conversion_component} gives the
transfer-potential evolution and component-compatibility relations.  Choosing
component \(N\) as the reference gives
%
\begin{empheq}[box=\fbox]{align}
	\frac{D_\xi\tau}{Dt}
	={} &
	\frac{1}{2}
	\mbf v_\xi\cdot\mbf v_\xi
	-
	\mu_\xi^N
	-
	\frac{\pi_\xi^N}{\phi_\xi\bar\rho_\xi\eta_\xi^N},
	\label{eq:MC_admissible_conversion_component}
	\\
	\mu_\xi^\alpha
	+
	\frac{\pi_\xi^\alpha}{\phi_\xi\bar\rho_\xi\eta_\xi^\alpha}
	={} &
	\mu_\xi^N
	+
	\frac{\pi_\xi^N}{\phi_\xi\bar\rho_\xi\eta_\xi^N},
	\qquad
	\alpha=1,\ldots,N-1.
	\label{eq:transfer_potential_component_compatibility}
\end{empheq}
%
Equation \eqref{eq:MC_admissible_conversion_component} evolves the single
Bedford--Drumheller transfer potential \(\tau\) using the reference component.
Equation \eqref{eq:transfer_potential_component_compatibility} then constrains
the remaining \(N-1\) component storage multipliers relative to component
\(N\).  Thus \(\pi_\xi^\alpha\) absorbs componentwise differences in the
thermodynamic potentials so that one transfer potential remains well-defined.

With this identification, substituting
\eqref{eq:reference_conversion_accumulation_rate} into
the temperature-weighted component-supply term in
\eqref{eq:MC_reduced_entropy_inequality} separates the reaction and
relative-transport contributions:
%
\begin{align}
	& -
	\sum_{\mc G\in\{\mc F,\mc S\}}
	\frac{1}{\theta_{\mc G}}
	\sum_{\xi\in\mc G}\sum_\alpha
	\phasedot{c}_\xi^\alpha
	\mu_\xi^\alpha
	\notag \\
	& \quad =
	-\sum_m
	\left[
		\sum_{\mc G\in\{\mc F,\mc S\}}
		\frac{1}{\theta_{\mc G}}
		\sum_{\xi\in\mc G}\sum_\alpha
		\mu_\xi^\alpha
		\nu_{\xi (m)}^\alpha
	\right]\dot r_{(m)}
	+
	\sum_{\mc G\in\{\mc F,\mc S\}}
	\sum_{\xi\in\mc G}\sum_\alpha
	\frac{\mu_\xi^\alpha}{\theta_{\mc G}}
	\nabla_{\mbf x}\cdot
		\left(
			\mbf j_{\xi,{\rm disp}}^\alpha
		+
			\mbf j_{\xi,{\rm diff}}^\alpha
		\right).
	\label{eq:MC_helmholtz_source_projection}
\end{align}
%
Here \(\mu_\xi^\alpha\) is the electrochemical potential defined by
\eqref{eq:C45_combined_def}.  Define the mechanism affinity using that same
potential:
%
\begin{equation}
	\mc A_{(m)}
	\equiv
	-
	\sum_\xi\sum_\alpha
	\mu_\xi^\alpha\nu_{\xi (m)}^\alpha .
	\label{eq:MC_affinity_projection}
\end{equation}
%
For a mechanism satisfying \eqref{eq:mechanism_charge_conservation},
substitution of
\(\mu_\xi^\alpha=\hat{\mu}_\xi^\alpha+z_\xi^\alpha\varphi\) gives
%
\begin{align}
	\mc A_{(m)}
	={} &
	-\sum_\xi\sum_\alpha
	\left(
	\hat{\mu}_\xi^\alpha+z_\xi^\alpha\varphi
	\right)
	\nu_{\xi (m)}^\alpha
	\notag \\
	={} &
	-\sum_\xi\sum_\alpha
	\hat{\mu}_\xi^\alpha\nu_{\xi (m)}^\alpha
	-
	\varphi
	\sum_\xi\sum_\alpha
	z_\xi^\alpha\nu_{\xi (m)}^\alpha
	\notag \\
	={} &
	-\sum_\xi\sum_\alpha
	\hat{\mu}_\xi^\alpha\nu_{\xi (m)}^\alpha .
	\label{eq:electrochemical_affinity_charge_cancellation}
\end{align}
%
Thus the explicit electric-potential contribution cancels from the affinity,
while the electric field still enters component transport, Gauss' law, and the
state dependence of the phase potentials.

Each mechanism uses one of the two existing subsystem temperatures,
%
\begin{equation}
	\theta_{(m)}
	=
	\begin{cases}
		\theta_{\mc F},
		& \text{for a fluid or fluid--solid interfacial mechanism},\\
		\theta_{\mc S},
		& \text{for a mechanism confined to the solid subsystem}.
	\end{cases}
	\label{eq:MC_reaction_temperature_assignment}
\end{equation}
%
If every nonzero stoichiometric coefficient belongs to one thermal subsystem,
the reaction term in \eqref{eq:MC_helmholtz_source_projection} is exactly
\(\mc A_{(m)}\dot r_{(m)}/\theta_{(m)}\).  For a fluid--solid mechanism, that
target affinity form requires an energy allocation between the separately
weighted phase terms.  The adopted interfacial temperature is
\(\theta_{(m)}=\theta_{\mc F}\).

\pagebreak[3]
The required allocation follows by subtracting the raw source contribution in
\eqref{eq:MC_helmholtz_source_projection} from the target affinity form,
%
\begin{align}
	& -
	\frac{\dot r_{(m)}}{\theta_{(m)}}
	\sum_{\mc G\in\{\mc F,\mc S\}}
	\sum_{\xi\in\mc G}\sum_\alpha
	\mu_\xi^\alpha\nu_{\xi (m)}^\alpha
	+
	\dot r_{(m)}
	\sum_{\mc G\in\{\mc F,\mc S\}}
	\frac{1}{\theta_{\mc G}}
	\sum_{\xi\in\mc G}\sum_\alpha
	\mu_\xi^\alpha\nu_{\xi (m)}^\alpha
	\notag \\
	& \quad =
	\sum_{\mc G\in\{\mc F,\mc S\}}
	\left(
	\frac{1}{\theta_{\mc G}}
	-
	\frac{1}{\theta_{(m)}}
	\right)
	\left(
	\sum_{\xi\in\mc G}\sum_\alpha
	\mu_\xi^\alpha\nu_{\xi (m)}^\alpha
	\right)
	\dot r_{(m)}
	\notag \\
	& \quad =
	\sum_{\mc G\in\{\mc F,\mc S\}}
	\frac{1}{\theta_{\mc G}}
	\left.
	\sum_{\xi\in\mc G}
	\rho_\xi e_\xi^{\rm int}
	\right\vert_{{\rm conv},(m)} .
	\label{eq:MC_reaction_energy_allocation}
\end{align}
%
Consequently, adding the allocation identified by
\eqref{eq:MC_reaction_energy_allocation} to the raw source contribution leaves
\(\mc A_{(m)}\dot r_{(m)}/\theta_{(m)}\).  Equation
\eqref{eq:MC_reaction_energy_allocation} is the two-temperature analogue of the
coupled conversion and heat-exchange terms in Drumheller's constituent entropy
inequality \cite[eq.~(88)]{drumheller2000}.

The last term is the relative-transport contribution.  Applying the product
rule to the temperature-weighted component potential gives
%
\begin{equation}
	\begin{aligned}
		 &
		\sum_{\mc G\in\{\mc F,\mc S\}}
		\sum_{\xi\in\mc G}\sum_\alpha
		\frac{\mu_\xi^\alpha}{\theta_{\mc G}}
		\nabla_{\mbf x}\cdot
		\left(
		\mbf j_{\xi,{\rm disp}}^\alpha
		+
		\mbf j_{\xi,{\rm diff}}^\alpha
		\right)
		\\
		 & \quad =
		\sum_{\mc G\in\{\mc F,\mc S\}}
		\sum_{\xi\in\mc G}\sum_\alpha
		\nabla_{\mbf x}\cdot
		\left[
			\frac{\mu_\xi^\alpha}{\theta_{\mc G}}
			\left(
			\mbf j_{\xi,{\rm disp}}^\alpha
			+
			\mbf j_{\xi,{\rm diff}}^\alpha
			\right)
		\right]
		\\
		 & \qquad -
		\sum_{\mc G\in\{\mc F,\mc S\}}
		\sum_{\xi\in\mc G}\sum_\alpha
		\left(
		\mbf j_{\xi,{\rm disp}}^\alpha
		+
		\mbf j_{\xi,{\rm diff}}^\alpha
		\right)\cdot
		\nabla_{\mbf x}
		\left(
			\frac{\mu_\xi^\alpha}{\theta_{\mc G}}
		\right).
	\end{aligned}
	\label{eq:relative_transport_product_rule}
\end{equation}
%
Integrating \eqref{eq:relative_transport_product_rule} over an arbitrary current
region and applying the divergence theorem converts the first term on its
right-hand side into entropy transport through the boundary.  The remaining
term identifies the local relative-transport contribution as
%
\begin{equation}
	-\sum_{\mc G\in\{\mc F,\mc S\}}
	\sum_{\xi\in\mc G}\sum_\alpha
	\left(
	\mbf j_{\xi,{\rm disp}}^\alpha
	+
	\mbf j_{\xi,{\rm diff}}^\alpha
	\right)\cdot
	\nabla_{\mbf x}
	\left(
	\frac{\mu_\xi^\alpha}{\theta_{\mc G}}
	\right).
	\label{eq:MC_relative_transport_power}
\end{equation}
%
This temperature-weighted potential gradient is the standard species-flux
force in nonisothermal mixture thermodynamics
\cite{drumheller2000,bothe2015continuum}.  It reduces to the energy-potential
gradient used in the Onsager subsection when the local temperature is uniform.

\subsubsection{Residual inequality}
\label{sec:MC_residual_inequality}

The plastic-deformation power in the residual inequality is driven directly by
the material stress \(\bs\sigma_s'\) identified in
\eqref{eq:MC_material_stress_power_identity}.  This stress is mapped through
each solid phase's distension and true-deformation gradients to its plastic-flow
configuration.
Each true-plastic flow is isochoric, so the deviatoric part of that mapped
stress is conjugate to the true-plastic rate.  Define \(\mbf S_s\) as
the deviatoric true-plastic driving stress,
%
\begin{equation}
	\mbf S_s
	\equiv
	(\bar{\mbf F}_s^e)^{-1}
	\mbf A_s^{-1}
	\left(
	(\bs\sigma_s')^T
	-\frac{1}{3}
	\operatorname{tr}\left(
	\bs\sigma_s'
	\right)\mbf I
	\right)
	\mbf A_s
	\bar{\mbf F}_s^e,
	\qquad
	s\in\mc S.
	\label{eq:true_plastic_driving_stress_definition}
\end{equation}
%
Because \(\mbf S_s\) is a similarity transformation of the deviatoric stress,
\(\operatorname{tr}\mbf S_s=0\).  The spherical part is therefore orthogonal to
the isochoric rate
\(\dot{\bar{\mbf F}}_s^p(\bar{\mbf F}_s^p)^{-1}\).  Equation
\eqref{eq:true_plastic_driving_stress_definition} is the phase-distension
mapping of Drumheller's deviatoric true-plastic driving stress
\cite[eqs.~(108)--(112)]{drumheller2000}.

Starting from \eqref{eq:MC_collected_entropy}, substitute the reversible
restrictions
\eqref{eq:MC_phase_pressure_definitions}, \eqref{eq:MC_solid_stress_definition},
\eqref{eq:MC_distension_force_definition},
\eqref{eq:MC_volume_fraction_restrictions},
\eqref{eq:MC_fluid_composition_projection},
\eqref{eq:MC_solid_composition_projection}, and
\eqref{eq:MC_entropy_definitions}.  Apply the source and transport reductions
\eqref{eq:MC_helmholtz_source_projection},
\eqref{eq:MC_affinity_projection}, and
\eqref{eq:MC_relative_transport_power}, including the interfacial
energy-allocation condition \eqref{eq:MC_reaction_energy_allocation}.  Finally,
\eqref{eq:gauss_law} removes the electric-potential-rate coefficient.  The
remaining constitutive terms give the residual inequality
%
\begin{equation}
	\begin{aligned}
		0\le{}\; &
		\sum_m
		\frac{\mc A_{(m)}\dot r_{(m)}}{\theta_{(m)}}
		-
		\sum_{\mc G\in\{\mc F,\mc S\}}
		\sum_{\xi\in\mc G}\sum_\alpha
		\left(
		\mbf j_{\xi,{\rm disp}}^\alpha
		+
		\mbf j_{\xi,{\rm diff}}^\alpha
		\right)\cdot
		\nabla_{\mbf x}
			\left(
			\frac{\mu_\xi^\alpha}{\theta_{\mc G}}
			\right)
			\\
			         & +
			\frac{1}{\theta_{\mc S}}
		\sum_{s\in\mc S}
		\bs\sigma_s':
		\mbf A_s^e
		\dot{\mbf A}_s^p
		(\mbf A_s^p)^{-1}
		(\mbf A_s^e)^{-1}
		\\
		         & +
		\frac{1}{\theta_{\mc S}}
		\sum_{s\in\mc S}
		\operatorname{tr}\left[
			\mbf S_s
			\dot{\bar{\mbf F}}_s^p
			(\bar{\mbf F}_s^p)^{-1}
			\right]
		\\
		         & +
		\sum_{\mc G\in\{\mc F,\mc S\}}
		\frac{1}{\theta_{\mc G}}
		\sum_{\xi\in\mc G}
		\rho_\xi e_\xi^{\rm int}
		-
		\sum_{\mc G\in\{\mc F,\mc S\}}
		\sum_{\xi\in\mc G}
		\frac{
		\mbf q_\xi\cdot\nabla_{\mbf x}\theta_{\mc G}
		}{
		\theta_{\mc G}^2
		}.
	\end{aligned}
	\label{eq:MC_residual_dissipation}
\end{equation}
%
\subsubsection{Reaction and relative-transport closures}
\label{sec:MC_reaction_transport_closures}

Inequality \eqref{eq:MC_residual_dissipation} is the constitutive admissibility
condition.  The quadratic Onsager potential in
\sref{onsager_reaction_transport} provides one near-equilibrium specialization
that satisfies it.  For uncoupled mechanisms, the reaction closure is
%
\begin{equation}
	\dot r_{(m)}
	=
	K_{(m)}\mc A_{(m)},
	\qquad
	K_{(m)}>0 .
	\label{eq:MC_onsager_reaction_rate}
\end{equation}
%
The corresponding stationarity relations for dispersion and diffusion are
%
\begin{align}
	\sum_\beta
	\left(\mbf D_{\xi,{\rm disp}}^{-1}\right)^{\alpha\beta}
	\mbf j_{\xi,{\rm disp}}^\beta
	+
	\theta_{\mc G}
	\nabla_{\mbf x}
	\left(
	\frac{\mu_\xi^\alpha}{\theta_{\mc G}}
	\right)
	 & =\mbf 0,
	\qquad \xi\in\mc G,
	\label{eq:MC_onsager_dispersion_stationarity}
	\\
	\sum_\beta
	\left(\mbf D_{\xi,{\rm diff}}^{-1}\right)^{\alpha\beta}
	\mbf j_{\xi,{\rm diff}}^\beta
	+
	\theta_{\mc G}
	\nabla_{\mbf x}
	\left(
	\frac{\mu_\xi^\alpha}{\theta_{\mc G}}
	\right)
	 & =\mbf 0,
	\qquad \xi\in\mc G.
	\label{eq:MC_onsager_diffusion_stationarity}
\end{align}
%
or, equivalently,
%
\begin{align}
	\mbf j_{\xi,{\rm disp}}^\alpha
	 & =-\sum_\beta
	\mbf D_{\xi,{\rm disp}}^{\alpha\beta}
	\theta_{\mc G}
	\nabla_{\mbf x}
	\left(
	\frac{\mu_\xi^\beta}{\theta_{\mc G}}
	\right),
	\qquad \xi\in\mc G,
	\label{eq:MC_dispersion_closure}
	\\
	\mbf j_{\xi,{\rm diff}}^\alpha
	 & =-\sum_\beta
	\mbf D_{\xi,{\rm diff}}^{\alpha\beta}
	\theta_{\mc G}
	\nabla_{\mbf x}
	\left(
	\frac{\mu_\xi^\beta}{\theta_{\mc G}}
	\right),
	\qquad \xi\in\mc G,
	\label{eq:MC_diffusion_closure}
	\\
	\sum_\alpha\mbf D_{\xi,k}^{\alpha\beta}
	 & =\mbf 0,
	\qquad
	\text{for every }\beta,
	\qquad
	k\in\{{\rm disp},{\rm diff}\},
	\notag \\
	\sum_\beta\mbf D_{\xi,k}^{\alpha\beta}
	 & =\mbf 0,
	\qquad
	\text{for every }\alpha,
	\qquad
	k\in\{{\rm disp},{\rm diff}\}.
	\label{eq:MC_relative_transport_closures}
\end{align}
%
The zero column sums enforce the relative-flux constraint, while the zero row
sums make the closure invariant under a common shift of all component
potentials within a phase.  Together they make each mobility singular on the
full component space.  Accordingly, \(\mbf D_{\xi,k}^{-1}\) in
\eqref{eq:MC_onsager_dispersion_stationarity}--\eqref{eq:MC_onsager_diffusion_stationarity}
denotes the inverse restricted to the zero-sum relative-flux subspace.
Substitution of
\eqref{eq:MC_onsager_reaction_rate} and
\eqref{eq:MC_relative_transport_closures} into the reaction and transport part
of \eqref{eq:MC_residual_dissipation} gives
%
\begin{align}
	 & \sum_m
	\frac{\mc A_{(m)}\dot r_{(m)}}{\theta_{(m)}}
	-
	\sum_{\mc G\in\{\mc F,\mc S\}}
	\sum_{\xi\in\mc G}\sum_\alpha
	\left(
	\mbf j_{\xi,{\rm disp}}^\alpha
	+
	\mbf j_{\xi,{\rm diff}}^\alpha
	\right)\cdot
	\nabla_{\mbf x}
	\left(
	\frac{\mu_\xi^\alpha}{\theta_{\mc G}}
	\right)
	\notag                           \\
	 & \qquad=
	\sum_m
	\frac{\dot r_{(m)}^2}{K_{(m)}\theta_{(m)}}
	\notag \\
	 & \qquad\quad +
	\sum_{\mc G\in\{\mc F,\mc S\}}
	\frac{1}{\theta_{\mc G}}
	\sum_{\xi\in\mc G}\sum_\alpha\sum_\beta
	\mbf j_{\xi,{\rm disp}}^\alpha
	\cdot
	\left(\mbf D_{\xi,{\rm disp}}^{-1}\right)^{\alpha\beta}
	\mbf j_{\xi,{\rm disp}}^\beta
	\notag \\
	 & \qquad\quad +
	\sum_{\mc G\in\{\mc F,\mc S\}}
	\frac{1}{\theta_{\mc G}}
	\sum_{\xi\in\mc G}\sum_\alpha\sum_\beta
	\mbf j_{\xi,{\rm diff}}^\alpha
	\cdot
	\left(\mbf D_{\xi,{\rm diff}}^{-1}\right)^{\alpha\beta}
	\mbf j_{\xi,{\rm diff}}^\beta
	\ge0.
	\label{eq:MC_onsager_reaction_transport_dissipation}
\end{align}
%
Thus the reaction and relative-transport terms are dissipative when each
\(\mbf D_{\xi,k}\) is symmetric positive semidefinite on the full component
space and positive definite on the relative-flux subspace satisfying
\eqref{eq:averaged_relative_flux_sum}.  Diffusion and dispersion may use
different mobilities, but both are driven here by the same
temperature-weighted electrochemical-potential gradient; the neutral limit is
obtained by setting \(z_\xi^\alpha=0\).  Dispersion may depend on resolved flow
and pore geometry~\cite{plumb1988dispersion,gray2009tcat_transport}.  Reaction,
diffusion, and dispersion may also be cross-coupled, provided their combined
force--rate form remains nonnegative~\cite{doi2011onsager,wang2021onsager}.

\subsubsection{Plastic-flow closures}
\label{sec:MC_plastic_flow_closures}

The two plastic powers in \eqref{eq:MC_residual_dissipation} admit the following
uncoupled associated closures, analogous to Drumheller's sufficient
construction~\cite{drumheller2000}:
%
\begin{subequations}
	\label{eq:MC_generic_plastic_flow_rules}
	\begin{align}
		\dot{\bar{\mbf F}}_s^p(\bar{\mbf F}_s^p)^{-1}
		 & =
		\Lambda_{\bar{\mbf F}_s}(\mc X_s)
		(\mbf S_s)^T,
		\label{eq:MC_generic_plastic_deformation_flow_rule}            \\
		\dot{\mbf A}_s^p(\mbf A_s^p)^{-1}
		 & =
		\Lambda_{\mbf A_s}(\mc X_s)
		(\mbf A_s^e)^T
		\bs\sigma_s'
		(\mbf A_s^e)^{-T},
		\label{eq:MC_generic_plastic_distension_flow_rule}             \\
		\Lambda_{\bar{\mbf F}_s}(\mc X_s)
		 & \ge0,
		\label{eq:MC_generic_plastic_deformation_mobility_nonnegative} \\
		\Lambda_{\mbf A_s}(\mc X_s)
		 & \ge0,
		\qquad
		s\in\mc S .
		\label{eq:MC_generic_plastic_distension_mobility_nonnegative}
	\end{align}
\end{subequations}
%
Because \(\operatorname{tr}\mbf S_s=0\), the first closure preserves
isochoric true-plastic flow.  The nonnegative scalar mobilities make both
plastic powers nonnegative separately.  Their dependence on the state set
\(\mc X_s\) permits composition-dependent mobilities.  Hardening, softening,
and other history dependence require additional internal variables and
evolution equations beyond the state set in \eqref{eq:MC_state_sets}.
Component-resolved mobilities for the additive specialization are given in
\eqref{eq:MC_component_resolved_plastic_flow_rules}.  Coupled or non-associated
plastic deformation--distension laws may replace
\eqref{eq:MC_generic_plastic_flow_rules}, but their combined conjugate power in
\eqref{eq:MC_residual_dissipation} must remain nonnegative.

\subsubsection{Interphase interaction and heat transport}
\label{sec:MC_interaction_heat_transport}

The remaining terms in \eqref{eq:MC_residual_dissipation} contain interphase
mechanical power, thermal exchange, and heat conduction.  Following Bedford and
Drumheller~\cite{drumheller1980thermomechanical}, conservation of total mixture
energy requires the internal interaction-energy supplies and interphase
momentum supplies to satisfy
%
\begin{equation}
	\sum_\xi
	\left(
	\rho_\xi e_\xi^{\rm int}
	+\mbf b_\xi\cdot\mbf v_\xi
	\right)
	=0 .
	\label{eq:MC_interaction_energy_conservation}
\end{equation}
%
The two subsystem balances contain equal-and-opposite reversible
volume-constraint powers, whose temperature-weighted coefficients vanish
through \eqref{eq:MC_volume_fraction_restrictions}.  Equation
\eqref{eq:MC_interaction_energy_conservation} therefore governs internal-energy
and momentum exchange, while \eqref{eq:MC_residual_dissipation} collects the
remaining dissipative processes.

Decompose the net interphase force into pairwise forces,
%
\begin{equation}
	\mbf b_\xi
	=
	\sum_{\zeta\ne\xi}
	\mbf b_{\zeta\xi},
	\qquad
	\mbf b_{\zeta\xi}
	=
	-\mbf b_{\xi\zeta},
	\label{eq:MC_pairwise_interaction_force}
\end{equation}
%
where \(\mbf b_{\zeta\xi}\) is the force exerted on phase \(\xi\) by phase
\(\zeta\).  The rate transferred from relative mechanical motion to internal
energy is
%
\begin{equation}
	-\sum_{\xi<\zeta}
	\mbf b_{\zeta\xi}\cdot
	\left(
	\mbf v_\xi-\mbf v_\zeta
	\right).
	\label{eq:MC_interaction_power_pairing}
\end{equation}
%
A pairwise quadratic construction gives a sufficient dissipative closure
\cite{foster2026reactingporous}.  Introduce symmetric positive-definite
interphase mobility tensors \(\mbf k_{\xi\zeta}\), a nonnegative fluid--solid
heat-transfer coefficient \(h_{\mc F\mc S}\), and symmetric positive
semidefinite heat-conductivity tensors \(\mbf K_\xi\):
%
\begin{subequations}
	\label{eq:MC_interaction_heat_closures}
	\begin{align}
		\mbf b_{\zeta\xi}
		 & =
		-\mbf k_{\xi\zeta}^{-1}
		(\mbf v_\xi-\mbf v_\zeta),
		\qquad
		\mbf k_{\xi\zeta}
		=
		\mbf k_{\zeta\xi}
		=
		\mbf k_{\xi\zeta}^T>0,
		\label{eq:MC_interphase_exchange_law} \\
		\left.
		\sum_{f\in\mc F}\rho_f e_f^{\rm int}
		\right\vert_{\rm heat}
		 & =
		h_{\mc F\mc S}
		\left(
		\theta_{\mc S}-\theta_{\mc F}
		\right)
		=
		-
		\left.
		\sum_{s\in\mc S}\rho_s e_s^{\rm int}
		\right\vert_{\rm heat},
		\qquad
		h_{\mc F\mc S}\ge0,
		\label{eq:MC_fluid_solid_heat_exchange_law} \\
		\mbf q_\xi
		 & =
		-\mbf K_\xi\nabla_{\mbf x}\theta_{\mc G},
		\qquad
		\xi\in\mc G,
		\qquad
		\mbf K_\xi=\mbf K_\xi^T\ge0 .
		\label{eq:MC_heat_flux_law}
	\end{align}
\end{subequations}
%
where \(\mbf y\cdot\mbf k_{\xi\zeta}\mbf y>0\) for every
\(\mbf y\ne\mbf 0\), \(\mbf y\cdot\mbf K_\xi\mbf y\ge0\), and
\(\theta_{\mc F},\theta_{\mc S}>0\).  For a phase pair, the force law converts
the lost relative mechanical power into heat.  Splitting that heat equally
between the two phases gives
%
\begin{align}
	& \frac{1}{2}
	\sum_{\xi<\zeta}
	\left(
	\frac{1}{\theta_\xi}
	+
	\frac{1}{\theta_\zeta}
	\right)
	\left(
	\mbf v_\xi-\mbf v_\zeta
	\right)
	\cdot
	\mbf k_{\xi\zeta}^{-1}
	\left(
	\mbf v_\xi-\mbf v_\zeta
	\right)
	\ge0 ,
	\label{eq:MC_interphase_exchange_positive}
\end{align}
%
which is the drag contribution to entropy production.  The fluid--solid heat
exchange contributes
%
\begin{equation}
	h_{\mc F\mc S}
	\left(
	\theta_{\mc S}-\theta_{\mc F}
	\right)
	\left(
	\frac{1}{\theta_{\mc F}}
	-
	\frac{1}{\theta_{\mc S}}
	\right)
	=
	h_{\mc F\mc S}
	\frac{
	\left(
	\theta_{\mc S}-\theta_{\mc F}
	\right)^2
	}{
	\theta_{\mc F}\theta_{\mc S}
	}
	\ge0 .
	\label{eq:MC_fluid_solid_heat_exchange_positive}
\end{equation}
%
The heat-conduction part becomes
%
\begin{equation}
	-\sum_{\mc G\in\{\mc F,\mc S\}}
	\sum_{\xi\in\mc G}
	\frac{
	\mbf q_\xi\cdot\nabla_{\mbf x}\theta_{\mc G}
	}{
	\theta_{\mc G}^2
	}
	=
	\sum_{\mc G\in\{\mc F,\mc S\}}
	\sum_{\xi\in\mc G}
	\frac{
		\nabla_{\mbf x}\theta_{\mc G}
		\cdot
		\mbf K_\xi
		\nabla_{\mbf x}\theta_{\mc G}
	}{
	\theta_{\mc G}^2
	}
	\ge0 .
	\label{eq:MC_heat_flux_positive}
\end{equation}
%
Equations \eqref{eq:MC_onsager_reaction_transport_dissipation},
\eqref{eq:MC_interphase_exchange_positive}, and
\eqref{eq:MC_fluid_solid_heat_exchange_positive}--\eqref{eq:MC_heat_flux_positive},
together with admissible plastic laws, provide nonnegative reaction,
relative-transport, drag, heat-exchange, and conduction contributions for
distinct positive temperatures \(\theta_{\mc F}\) and \(\theta_{\mc S}\).

\subsection{Sorption/desorption mechanisms}
\label{sec:MC_sorption_desorption_mechanisms}

Sorption/desorption is a conversion mechanism that transfers a component
between a mobile phase and a solid phase while the solid constituents share the
skeleton motion.  For a component \(g\) moving between mobile phase \(f\) and solid phase
\(s\), the stoichiometric coefficients may be written
%
\begin{equation}
	\nu_{s (m)}^g=1,
	\qquad
	\nu_{f (m)}^g=-1,
	\qquad
	\nu_{\xi (m)}^\alpha=0
	\quad\hbox{otherwise}.
	\label{eq:MC_sorption_stoichiometry}
\end{equation}
%
The affinity is
%
\begin{equation}
	\mc A_{(m)}
	=
	\mu_f^g-\mu_s^g .
	\label{eq:MC_sorption_affinity}
\end{equation}
%
An admissible kinetic law must satisfy
\(\mc A_{(m)}\dot r_{(m)}\ge0\).  The sorbed amount \(\eta_s^g\) can also
enter \(\psi_s\), \(\mbf A_s^0\), \(\bar{\mbf F}_s^0\), and the
plastic flow law, so sorption can change stiffness, yield strength, swelling,
and compaction behavior in one thermodynamic state.

Hydrate formation and dissociation can be represented in the same
sorption/desorption form.  For methane hydrate or another clathrate solid, host
water and guest methane are solid components of the same load-bearing solid
phase, while exchange with mobile components changes the solid composition.
That composition may change the free energy, thermal expansion, stress-free
cage volume, and plastic response within the shared solid displacement field.

\subsection{Phase reaction-modified Biot coefficients}
\label{sec:MC_phase_biot_legendre_transform}

The restrictions in \sref{MC_reversible_restrictions} define the single-prime
solid stress directly from the phase Helmholtz energy.  That stress is paired
with the elastic rates in the entropy inequality when the intrinsic solid
density is fixed.  The corresponding fixed-equivalent-pressure double-prime
stress is obtained by a phase-level Legendre transform in the standard
thermodynamic sense \cite{callen1985thermodynamics}, before any additive
component specialization is introduced.

Define the phase and intrinsic specific volumes
%
\begin{equation}
	v_s=\frac{1}{\rho_s},
	\qquad
	v_{s0}=\frac{1}{\rho_{s0}},
	\qquad
	\bar v_s=\frac{1}{\bar\rho_s},
	\qquad
	\bar v_{s0}=\frac{1}{\bar\rho_{s0}},
	\qquad
	s\in\mc S .
	\label{eq:MC_phase_specific_volumes}
\end{equation}
%
The solid volume-fraction restriction and the phase pressure definition give
the equivalent pressure in specific-volume form,
%
\begin{align}
	\bar p_E
	 & =
	\bar p_s-\omega_s^+
	\notag \\
	 & =
	-
	\left.
	\p{
		\left(
		\psi_s+\omega_s^+\bar v_s
		\right)
	}{\bar v_s}
	\right\vert_{
		\bar{\mbf F}_s^e,\mbf A_s^e,
		\{\eta_s^\alpha\},\theta_{\mc S},\mbf E
	},
	\qquad
	s\in\mc S ,
	\label{eq:MC_phase_pressure_specific_volume}
\end{align}
%
where \(\bar v_s\) is the specific-volume representation of the density
argument \(\bar\rho_s\).  Thus the electric enthalpy per intrinsic solid mass
is \(\omega_s^+\bar v_s\), and the phase Legendre transform is
%
\begin{subequations}
	\label{eq:MC_phase_legendre_transform}
	\begin{gather}
		\widetilde\psi_s
		=
		\psi_s
		+\omega_s^+\bar v_s
		+\bar p_E\bar v_s,
		\label{eq:MC_phase_legendre_transform_definition} \\
		\bar v_s
		=
		\left.
		\p{\widetilde\psi_s}{\bar p_E}
		\right\vert_{\bar{\mbf F}_s^e,\mbf A_s^e,
			\{\eta_s^\alpha\},\theta_{\mc S},\mbf E}.
		\label{eq:MC_phase_legendre_transform_pressure_derivative}
	\end{gather}
\end{subequations}
%
For the fixed-pressure elastic derivative, the density dependence induced by
\(\bar{\mbf F}_s^e\) enters through the true Jacobian \(J=a_s\bar J_s\).  This
follows from the material mass balance \eqref{eq:material_mass_2}, and from its
true-solid specialization \eqref{eq:MC_solid_true_mass_conservation}: density
and specific volume depend on the local Jacobian.  The current values of the
distension, plastic, stress-free, thermal, and composition variables are part of
the state at which the tangent is evaluated, while the differentiation is taken
at fixed equivalent pressure, electric field, solid temperature, and composition,
%
\begin{align}
	\left.
	\p{\bar v_s}{\bar{\mbf F}_s^e}
	\right\vert_{\bar p_E,\theta_{\mc S},\{\eta_s^\alpha\},\mbf E}
	 & =
	\left.
	\p{\bar v_s}{J}
	\right\vert_{\bar p_E,\theta_{\mc S},\{\eta_s^\alpha\},\mbf E}
	\frac{\partial J}{\partial\bar{\mbf F}_s^e}
	\notag \\
	 & =
	\left.
	\p{\bar v_s}{J}
	\right\vert_{\bar p_E,\theta_{\mc S},\{\eta_s^\alpha\},\mbf E}
	a_s\bar J_s
	(\bar{\mbf F}_s^e)^{-T}.
	\label{eq:MC_phase_specific_volume_deformation_chain_rule}
\end{align}
%
The phase electric enthalpy per intrinsic solid mass is
\(\omega_s^+\bar v_s\).  The single-prime and double-prime phase Piola stresses
are defined by
%
\begin{subequations}
	\label{eq:MC_phase_pk_stresses}
	\begin{align}
		\mbf P_s'
		 & =
		\rho_{s0}
		\left.
		\p{\widetilde\psi_s}{\bar{\mbf F}_s^e}
		\right\vert_{\bar p_E,\theta_{\mc S},\{\eta_s^\alpha\},\mbf E}
		\label{eq:MC_phase_single_prime_pk_stress} \\
		\rho_{s0}
		\left.
		\p{\widetilde\psi_s}{\bar{\mbf F}_s^e}
		\right\vert_{\bar p_E,\theta_{\mc S},\{\eta_s^\alpha\},\mbf E}
		 & =
		\rho_{s0}
		\left.
		\p{\psi_s}{\bar{\mbf F}_s^e}
		\right\vert_{\bar v_s,\{\eta_s^\alpha\},\theta_{\mc S},\mbf E},
		\label{eq:MC_phase_legendre_pk_identity}   \\
		\mbf P_s^{\prime\prime}
		 & =
		\rho_{s0}
		\left.
		\p{
			\left(
			\psi_s+\omega_s^+\bar v_s
			\right)
		}{\bar{\mbf F}_s^e}
		\right\vert_{\bar p_E,\theta_{\mc S},\{\eta_s^\alpha\},\mbf E}.
		\label{eq:MC_phase_double_prime_pk_stress}
	\end{align}
\end{subequations}
%
Using the Legendre transform at fixed \(\bar p_E\),
%
\begin{equation}
	\mbf P_s'
	=
	\mbf P_s^{\prime\prime}
	+
	\bar p_E\rho_{s0}
	\left.
	\p{\bar v_s}{\bar{\mbf F}_s^e}
	\right\vert_{\bar p_E,\theta_{\mc S},\{\eta_s^\alpha\},\mbf E}.
	\label{eq:MC_phase_pk_legendre_split}
\end{equation}
%
Substitution of \eqref{eq:MC_phase_specific_volume_deformation_chain_rule}
identifies the phase nonlinear Biot coefficient
%
\begin{empheq}[box=\fbox]{equation}
	B_s
	\equiv
	1
	-
	\frac{a_s}{v_{s0}}
	\left.
	\p{\bar v_s}{J}
	\right\vert_{\bar p_E,\theta_{\mc S},\{\eta_s^\alpha\},\mbf E},
	\qquad
	s\in\mc S .
	\label{eq:MC_phase_biot_coefficient}
\end{empheq}
%
The coefficient \(B_s\) is a skeleton-volume tangent evaluated in the current
reacted state.  The factor \(a_s\) maps the common skeleton Jacobian to the true
solid Jacobian and therefore carries the current pore-space allocation, packing,
and stress-free volume accommodation into the coefficient.  Thus the Biot
coefficient is derived from the reacted solid state and the fixed-pressure
stress split.  When the distension split is
suppressed, this reduces to the nonlinear Biot coefficient form derived in
\cite{foster2025}.  The phase Piola split is
%
\begin{equation}
	\mbf P_s'
	=
	\mbf P_s^{\prime\prime}
	+
	\left(1-B_s\right)
	\bar p_E
	\bar J_s
	(\bar{\mbf F}_s^e)^{-T}.
	\label{eq:MC_phase_pk_biot_split}
\end{equation}
%
Pushing the Piola stresses forward by
\(\bar J_s^{-1}(\cdot)(\bar{\mbf F}_s^e)^T\), and then through the same
distension map used in \eqref{eq:MC_solid_stress_definition}, leaves the
pressure correction isotropic.  Hence the phase Cauchy stress split is
%
\begin{empheq}[box=\fbox]{equation}
	\bs\sigma_s'
	=
	\bs\sigma_s^{\prime\prime}
	+
	\left(1-B_s\right)\bar p_E\mbf I,
	\qquad
	s\in\mc S .
	\label{eq:MC_phase_biot_split}
\end{empheq}
%
This phase relation applies to the general composition-dependent solid
Helmholtz energy.  The additive model in
\sref{MC_additive_component_specialization} supplies component formulas for
\(B_s\) and \(\bs\sigma_s^{\prime\prime}\).

\subsection{Phase and overall momentum after the Biot transform}
\label{sec:MC_phase_overall_momentum_biot}

Using the equivalent pressure \(\bar p_E=-\lambda\), the solid
volume-fraction restriction \eqref{eq:MC_solid_volume_fraction_restriction}
gives
%
\begin{align}
	\bar p_E
	 & =
	-\lambda
	\notag \\
	 & =
	\bar p_s-\omega_s^+,
	\qquad
	s\in\mc S .
	\label{eq:MC_phase_equivalent_pressure}
\end{align}
%
Substituting \eqref{eq:MC_phase_biot_split} into the solid phase momentum
balance \eqref{eq:el_momentum_uniform} gives the fixed-pressure form
%
\begin{equation}
	\begin{aligned}
		\rho_s\mbf a_s
		={} &
		\nabla_{\mbf x}\cdot
		\left[
			\bs\sigma_s^{\prime\prime}
			+
			\left(1-B_s\right)\bar p_E\mbf I
			+
			\bs\sigma_s^+
			\right]
		-\phi_s\nabla_{\mbf x}\bar p_E
		-\nabla_{\mbf x}\left(\phi_s\omega_s^+\right)
		+\rho_s\mbf g+\mbf b_s
		\\
		    & +
		\sum_\alpha\phasedot{c}_s^\alpha
		\left(
		\nabla_{\mbf x}\tau-\mbf v_s
		\right),
		\qquad
		s\in\mc S .
	\end{aligned}
	\label{eq:MC_solid_phase_momentum_biot}
\end{equation}
%
The plasticity restrictions are written in terms of
\(\bs\sigma_s'\), whose conjugate power remains after the electric-enthalpy
rate accounts for the Maxwell power.  If a plastic model is calibrated in terms of
\(\bs\sigma_s^{\prime\prime}\), the phase split
\eqref{eq:MC_phase_biot_split} must be applied before evaluating the plastic
driving stress.  The plastic flow law depends on \(\bs\sigma_s'\) and
composition; Maxwell stress affects that material stress through the coupled
momentum solution.

After the solid partial stresses are summed, the solid-phase Biot terms combine
as
%
\begin{align}
	\sum_{s\in\mc S}\bs\sigma_s'
	 & =
	\sum_{s\in\mc S}\bs\sigma_s^{\prime\prime}
	+
	\sum_{s\in\mc S}
	\left(1-B_s\right)\bar p_E\mbf I,
	\label{eq:MC_solid_biot_rollup_expanded}   \\
	\bs\sigma^{\prime\prime}
	 & \equiv
	\sum_{s\in\mc S}\bs\sigma_s^{\prime\prime},
	\label{eq:MC_solid_biot_rollup_definition} \\
	\Aboxed{
		B
	 & =
		1-\sum_{s\in\mc S}\left(1-B_s\right)
	}
	.
	\label{eq:MC_solid_biot_rollup_coefficient}
\end{align}
%
Thus the overall Biot coefficient is found after the pressure terms from the
solid phases have been added.  Each \(B_s\) is defined from the phase partial
Piola stress per unit skeleton reference volume, so the roll-up is the direct
sum in \eqref{eq:MC_solid_biot_rollup_coefficient}.  In the single-solid,
single-fluid limit, \(B=B_s\).

After summing the solid phase balances, the Biot-pressure and electrical parts in
\eqref{eq:MC_solid_phase_momentum_biot} can be rewritten using
\eqref{eq:MC_solid_biot_rollup_definition},
%
\begin{equation}
	\begin{aligned}
		 & \sum_{s\in\mc S}
		\left[
			\nabla_{\mbf x}\cdot
			\left(
			\left(1-B_s\right)\bar p_E\mbf I
			+\bs\sigma_s^+
			\right)
			-\phi_s\nabla_{\mbf x}\bar p_E
			-\nabla_{\mbf x}\left(\phi_s\omega_s^+\right)
			\right]
		\\
		 & \qquad =
		\nabla_{\mbf x}\cdot
		\left[
			-B\bar p_E\mbf I
			+
			\sum_{s\in\mc S}
			\phi_s\mbf E\otimes\mbf d_s
			\right]
		+
		\left(
		1-\sum_{s\in\mc S}\phi_s
		\right)
		\nabla_{\mbf x}\bar p_E .
	\end{aligned}
	\label{eq:solid_biot_pressure_remainder_rollup}
\end{equation}
%
The volume-fraction constraint identifies the coefficient of
\(\nabla_{\mbf x}\bar p_E\) in this solid remainder with the total fluid volume
fraction \(\phi=\sum_{f\in\mc F}\phi_f\).

For the fluid phases, substitute
\(\bar p_f=\bar p_E+\gamma_f+\omega_f^+\) and
\(\bs\sigma_f^+=\phi_f(\omega_f^+\mbf I+\mbf E\otimes\mbf d_f)\).
Writing \(\phi_f=\phi S_f\) and using the capillary identity
\(\sum_fS_f\nabla_{\mbf x}\gamma_f=\mbf0\) gives
%
\begin{equation}
	\sum_{f\in\mc F}
	\left[
		\nabla_{\mbf x}\cdot\bs\sigma_f^+
		-\phi_f\nabla_{\mbf x}\bar p_f
		-\omega_f^+\nabla_{\mbf x}\phi_f
		\right]
	=
	\nabla_{\mbf x}\cdot
	\left(
	\sum_{f\in\mc F}
	\phi_f\mbf E\otimes\mbf d_f
	\right)
	-\phi\nabla_{\mbf x}\bar p_E .
	\label{eq:MC_fluid_pressure_gradient_rollup}
\end{equation}
%
Adding \eqref{eq:solid_biot_pressure_remainder_rollup} and
\eqref{eq:MC_fluid_pressure_gradient_rollup}, using
\(\sum_\xi\phi_\xi=1\) and
\(\mbf d=\sum_\xi\phi_\xi\mbf d_\xi\), gives
%
\begin{equation}
	\begin{aligned}
		 & \sum_{s\in\mc S}
		\left\{
		\nabla_{\mbf x}\cdot
		\left[
			\left(1-B_s\right)\bar p_E\mbf I
			+\bs\sigma_s^+
			\right]
		-\phi_s\nabla_{\mbf x}\bar p_E
		-\nabla_{\mbf x}\left(\phi_s\omega_s^+\right)
		\right\}
		\\
		 & \quad
		+\sum_{f\in\mc F}
		\left[
			\nabla_{\mbf x}\cdot\bs\sigma_f^+
			-\phi_f\nabla_{\mbf x}\bar p_f
			-\omega_f^+\nabla_{\mbf x}\phi_f
			\right]
		\\
		 & \qquad =
		\nabla_{\mbf x}\cdot
		\left[
			-B\bar p_E\mbf I
			+\mbf E\otimes\mbf d
			\right].
	\end{aligned}
	\label{eq:MC_pressure_gradient_overall_rollup}
\end{equation}
%
Summing the phase momentum balances, using the pairwise internal momentum
antisymmetry \(\mbf b_{\zeta\xi}=-\mbf b_{\xi\zeta}\) in
\eqref{eq:MC_pairwise_interaction_force}, so that
\(\sum_\xi\sum_{\zeta\ne\xi}\mbf b_{\zeta\xi}=\mbf 0\), and substituting
\eqref{eq:MC_pressure_gradient_overall_rollup} gives the compact overall balance
%
\begin{empheq}[box=\fbox]{equation}
	\sum_\xi\rho_\xi\mbf a_\xi
	=
	\nabla_{\mbf x}\cdot
	\left[
		\bs\sigma^{\prime\prime}
		-
		B\bar p_E\mbf I
		+
		\mbf E\otimes\mbf d
		\right]
	+\rho\mbf g
	+\sum_\xi\sum_\alpha\phasedot{c}_\xi^\alpha
	\left(
	\nabla_{\mbf x}\tau-\mbf v_\xi
	\right).
	\label{eq:MC_overall_momentum_nonlinear_biot}
\end{empheq}

\subsection{Additive solid component specialization}
\label{sec:MC_additive_component_specialization}

The phase-level derivation allows a general solid Helmholtz energy.  The
additive model is useful when a solid phase has one shared skeleton motion but
its stored energy can still be split into component contributions.  Examples
include adsorbed species in a
mineral framework, bound water, substitutions or defects in a mineral lattice,
and guest occupancy in hydrate cages.  In these cases the components share the
same solid velocity and the same deformation measures
\(\bar{\mbf F}_s^e\) and \(\mbf A_s^e\).  The additive split is therefore an
energy representation for a common solid motion.  Component potentials may
depend on the full composition \(\{\eta_s^\alpha\}\), so nonideal component
interactions enter through the composition dependence of \(\psi_s^\alpha\).
The generic nonadditive formulation \eqref{eq:MC_solid_general_free_energy}
remains available when a component-energy split is too restrictive.
The component-level fixed-pressure formulas below use the neutral-solid
specialization \(\omega_s^+=0\), for which
\(\bar p_E=\bar p_s\).  The phase-level electrical response remains governed by
\eqref{eq:MC_phase_pressure_specific_volume}--\eqref{eq:MC_overall_momentum_nonlinear_biot}.

The additive model is obtained by choosing component potentials
\(\psi_s^\alpha\), component intrinsic densities \(\bar\rho_s^\alpha\), and
component pressures \(\bar p_s^\alpha\), then writing
%
\begin{equation}
	\psi_s
	=
	\sum_\alpha
	\eta_s^\alpha
	\psi_s^\alpha(
	\bar{\mbf F}_s^e,\mbf A_s^e,
	\bar\rho_s^\alpha,\{\eta_s^\beta\},\theta_{\mc S}).
	\label{eq:MC_additive_free_energy}
\end{equation}
%
Then the pressure, stress, and distension restrictions become
%
\begin{subequations}
	\label{eq:MC_additive_mixture_rules}
	\begin{align}
		\bar p_s
		 & =
		\sum_\alpha
		\phi_s^\alpha\bar p_s^\alpha,
		\label{eq:MC_additive_pressure_average}   \\
		\bar p_s^\alpha
		 & =
		(\bar\rho_s^\alpha)^2
		\left.
		\p{\psi_s^\alpha}{\bar\rho_s^\alpha}
		\right\vert_{\mc X_s},
		\label{eq:MC_additive_component_pressure} \\
		\bs\sigma_s'
		 & =
		\sum_\alpha
		\bs\sigma_s^{\prime\alpha},
		\label{eq:MC_additive_stress_sum}         \\
		\bs\sigma_s^{\prime\alpha}
		 & =
		\rho_s^\alpha
		\mbf A_s^{-T}
		\left(
		\left.
		\p{\psi_s^\alpha}{\bar{\mbf F}_s^e}
		\right\vert_{\mc X_s}
		(\bar{\mbf F}_s^e)^T
		\right)
		\mbf A_s^T,
		\label{eq:MC_additive_component_stress}   \\
		\bs\sigma_s'
		 & =
		\rho_s
		\sum_\alpha
		\eta_s^\alpha
		\left.
		\p{\psi_s^\alpha}{\mbf A_s^e}
		\right\vert_{\mc X_s}
		(\mbf A_s^e)^T .
		\label{eq:MC_additive_distension_force}
	\end{align}
\end{subequations}
%
The component pressures in
\eqref{eq:MC_additive_component_pressure} can also be read from the storage
multipliers.  The phase pressure-storage relation
\eqref{eq:phase_pressure_storage} gives
\(\sum_\alpha\pi_s^\alpha/\phi_s=\bar p_s\).  In the additive model, each term
in that sum is the component's share of the phase-pressure average,
%
\begin{equation}
	\frac{\pi_s^\alpha}{\phi_s}
	=
	\phi_s^\alpha\bar p_s^\alpha,
	\qquad
	\bar p_s^\alpha
	=
	\frac{\pi_s^\alpha}{\phi_s\phi_s^\alpha},
	\qquad
	s\in\mc S .
	\label{eq:additive_component_pressure_multiplier_identity}
\end{equation}
%
Summing this identity over components recovers
\eqref{eq:MC_additive_pressure_average} and
\eqref{eq:phase_pressure_storage}.  Thus the component pressures are
determined once the storage multipliers and component solid-volume
fractions are known.

The residual inequality uses the phase plastic powers.  A conservative
component-resolved plastic law therefore lets each component change the mobility
of the common phase plastic rates.  One sufficient additive form is
%
\begin{subequations}
	\label{eq:MC_component_resolved_plastic_flow_rules}
	\begin{align}
		\dot{\bar{\mbf F}}_s^p(\bar{\mbf F}_s^p)^{-1}
		 & =
		\left(
		\sum_\alpha
		\Lambda_{\bar{\mbf F}_s}^\alpha(\mc X_s)
		\right)
		(\mbf S_s)^T,
		\label{eq:MC_component_true_plastic_flow_rule}       \\
		\dot{\mbf A}_s^p(\mbf A_s^p)^{-1}
		 & =
		\left(
		\sum_\alpha
		\Lambda_{\mbf A_s}^\alpha(\mc X_s)
		\right)
		(\mbf A_s^e)^T
		\bs\sigma_s'
		(\mbf A_s^e)^{-T},
		\label{eq:MC_component_distension_plastic_flow_rule} \\
		\sum_\alpha
		\Lambda_{\bar{\mbf F}_s}^\alpha(\mc X_s)
		 & \ge0,
		\qquad
		\sum_\alpha
		\Lambda_{\mbf A_s}^\alpha(\mc X_s)
		\ge0 .
		\label{eq:MC_component_plastic_mobilities}
	\end{align}
\end{subequations}
%
These rules assign separate component mobility contributions to the common solid
plastic rates.  The summed mobilities enter the phase plastic powers and are
nonnegative.  Componentwise nonnegative mobilities make every component
contribution separately dissipative.
More general flow directions may depend on component identity, but the resulting
combined plastic power in \eqref{eq:MC_residual_dissipation} must remain
nonnegative.

\subsubsection{Component reaction-modified Biot coefficient derivation}
\label{sec:MC_component_biot_derivation}

The additive rules \eqref{eq:MC_additive_mixture_rules} give the component
single-prime stress.  The pressure-based momentum balance needs the
double-prime stress, so the fixed-pressure Legendre transform is applied to each
component before the component stresses are summed.
%
Let \(\bar v_s^\alpha=1/\bar\rho_s^\alpha\),
\(v_{s0}^\alpha=1/\rho_{s0}^\alpha\), and
\(\widetilde\psi_s^\alpha
=\psi_s^\alpha+\bar p_s^\alpha\bar v_s^\alpha\).  For each component, the
fixed-pressure Legendre transform follows the phase calculation in
\S\ref{sec:MC_phase_biot_legendre_transform}, with \(\bar p_s^\alpha\) replacing
\(\bar p_s\) and \(\bar v_s^\alpha\) replacing \(\bar v_s\).  The component
reference energy per unit skeleton reference volume is
\(\rho_{s0}^\alpha\psi_s^\alpha\), so the single-prime and
double-prime component partial Piola stresses are
%
\begin{subequations}
	\label{eq:MC_additive_component_pk_stresses}
	\begin{align}
		\mbf P_s^{\prime\alpha}
		 & =
		\rho_{s0}^\alpha
		\left.
		\p{\widetilde\psi_s^\alpha}{\bar{\mbf F}_s^e}
		\right\vert_{\bar p_s^\alpha,\theta_{\mc S},\{\eta_s^\beta\}}
		\label{eq:MC_additive_component_single_prime_pk_stress} \\
		\rho_{s0}^\alpha
		\left.
		\p{\widetilde\psi_s^\alpha}{\bar{\mbf F}_s^e}
		\right\vert_{\bar p_s^\alpha,\theta_{\mc S},\{\eta_s^\beta\}}
		 & =
		\rho_{s0}^\alpha
		\left.
		\p{\psi_s^\alpha}{\bar{\mbf F}_s^e}
		\right\vert_{\bar v_s^\alpha,\{\eta_s^\beta\},\theta_{\mc S}},
		\label{eq:MC_additive_component_legendre_pk_identity}   \\
		\mbf P_s^{\prime\prime\alpha}
		 & =
		\rho_{s0}^\alpha
		\left.
		\p{\psi_s^\alpha}{\bar{\mbf F}_s^e}
		\right\vert_{\bar p_s^\alpha,\theta_{\mc S},\{\eta_s^\beta\}}.
		\label{eq:MC_additive_component_double_prime_pk_stress}
	\end{align}
\end{subequations}
%
Using the Legendre transform at fixed \(\bar p_s^\alpha\),
%
\begin{equation}
	\mbf P_s^{\prime\alpha}
	=
	\mbf P_s^{\prime\prime\alpha}
	+
	\bar p_s^\alpha\rho_{s0}^\alpha
	\left.
	\p{\bar v_s^\alpha}{\bar{\mbf F}_s^e}
	\right\vert_{\bar p_s^\alpha,\theta_{\mc S},\{\eta_s^\beta\}}.
	\label{eq:MC_additive_component_pk_legendre_split}
\end{equation}
%
Using the same \(J=a_s\bar J_s\) chain rule as in the phase-level derivation,
the pressure correction in \eqref{eq:MC_additive_component_pk_legendre_split}
identifies the component nonlinear Biot coefficient
%
\begin{empheq}[box=\fbox]{equation}
	B_s^\alpha
	\equiv
	1
	-
	\frac{a_s}{v_{s0}^\alpha}
	\left.
	\p{\bar v_s^\alpha}{J}
	\right\vert_{\bar p_s^\alpha,\theta_{\mc S},\{\eta_s^\beta\}}.
	\label{eq:MC_additive_component_biot}
\end{empheq}
%
where \(B_s^\alpha\) is the component nonlinear Biot coefficient.  This gives the
component Piola split
%
\begin{equation}
	\mbf P_s^{\prime\alpha}
	=
	\mbf P_s^{\prime\prime\alpha}
	+
	(1-B_s^\alpha)\bar p_s^\alpha
	\bar J_s
	(\bar{\mbf F}_s^e)^{-T}.
	\label{eq:MC_additive_component_pk_biot_split}
\end{equation}
%
Pushing either component Piola stress forward by the operation
\(\bar J_s^{-1}(\cdot)(\bar{\mbf F}_s^e)^T\), and then through the same
distension map used in \eqref{eq:MC_additive_component_stress}, leaves the
pressure correction isotropic.  Hence the component Cauchy stress split is
%
\begin{equation}
	\bs\sigma_s^{\prime\alpha}
	=
	\bs\sigma_s^{\prime\prime\alpha}
	+
	(1-B_s^\alpha)\bar p_s^\alpha\mbf I .
	\label{eq:MC_additive_component_biot_split}
\end{equation}
%
Summing over components gives the phase stress split
%
\begin{subequations}
	\label{eq:MC_additive_biot_mixture_rule}
	\begin{align}
		\bs\sigma_s^{\prime\prime}
		 & =
		\sum_\alpha
		\bs\sigma_s^{\prime\prime\alpha},
		\label{eq:MC_additive_double_prime_sum} \\
		\bs\sigma_s'
		 & =
		\bs\sigma_s^{\prime\prime}
		+
		(1-B_s)\bar p_s\mbf I,
		\label{eq:MC_additive_phase_biot_split}
	\end{align}
	\begin{empheq}[box=\fbox]{equation}
		B_s
		=
		1
		-
		\frac{1}{\bar p_s}
		\sum_\alpha
		(1-B_s^\alpha)\bar p_s^\alpha.
		\label{eq:MC_additive_biot_mixture_coefficient}
	\end{empheq}
\end{subequations}
%
These additive rules give component formulas for the same phase stress split.
Equation \eqref{eq:MC_additive_biot_mixture_coefficient} combines the component
pressure corrections obtained by applying the fixed-pressure Legendre transform
to each component energy.  Thus the additive specialization shows how component
storage, composition dependence, and stress-free volume accommodation determine
the phase Biot coefficient, while retaining the plastic, distension, thermal,
and composition-dependent stress-free terms from the full entropy inequality.

\subsection{Single-component solid specialization}
\label{sec:MC_single_component_rollup}

The single-component mineral solid is recovered from the additive theory by
choosing one active component in each solid phase and imposing
%
\begin{equation}
	\eta_s^s=1,
	\qquad
	\dot\eta_s^s=0,
	\qquad
	\phi_s^s=1,
	\qquad
	\bar\rho_s^s=\bar\rho_s .
	\label{eq:MC_single_component_constraints}
\end{equation}
%
Then \eqref{eq:MC_solid_general_free_energy} reduces to the one-component solid
free energy
%
\begin{equation}
	\psi_s
	=
	\psi_s(
	\bar{\mbf F}_s^e,\mbf A_s^e,\bar\rho_s,\theta_{\mc S}),
	\label{eq:MC_single_component_free_energy}
\end{equation}
%
the solid composition projection \eqref{eq:MC_solid_composition_projection}
vanishes identically, and the additive component rules reduce to
%
\begin{equation}
	\bar p_s^s=\bar p_s,
	\qquad
	\bs\sigma_s^{\prime s}=\bs\sigma_s',
	\qquad
	\bs\sigma_s^{\prime\prime s}=\bs\sigma_s^{\prime\prime},
	\qquad
	B_s^s=B_s .
	\label{eq:MC_single_component_reductions}
\end{equation}
%
This identifies the one-component solid as the one-component limit of the
multicomponent solid.  The stress-free thermal maps, plastic variables,
distension variables, and source terms still belong to the phase, while the
solid-composition variation collapses to \(\dot\eta_s^s=0\).

For this specialization, the phase-level Biot coefficient
\eqref{eq:MC_phase_biot_coefficient}, the phase stress split
\eqref{eq:MC_phase_biot_split}, the solid phase momentum balance
\eqref{eq:MC_solid_phase_momentum_biot}, and the overall balance
\eqref{eq:MC_overall_momentum_nonlinear_biot} reduce to the one-component
solid formulas.  In the single-solid, single-fluid limit, the overall mixture
coefficient
\eqref{eq:MC_solid_biot_rollup_coefficient} gives \(B=B_s\).

\subsection{Scalar distension simplifications}
\label{sec:MC_scalar_distension_simplifications}

The tensor distension theory reduces to a scalar factor for isotropic
pore-structure change.  Scalar distension represents local solid-volume change,
while \(\mbf A_s\) represents anisotropic pore collapse and fabric evolution.  The scalar
specialization is imposed on the elastic, plastic, and stress-free distension
factors,
%
\begin{equation}
	\mbf A_s^\bullet=(a_s^\bullet)^{1/3}\mbf I,
	\qquad
	\bullet\in\{e,p,0\},
	\qquad
	s\in\mc S .
	\label{eq:MC_scalar_distension_factors}
\end{equation}
%
The factors multiply,
%
\begin{equation}
	a_s=a_s^e a_s^p a_s^0,
	\qquad
	s\in\mc S ,
	\label{eq:MC_scalar_distension_product}
\end{equation}
%
and the true solid deformation relation becomes
%
\begin{equation}
	\mbf F=a_s^{1/3}\bar{\mbf F}_s,
	\qquad
	J=a_s\bar J_s .
	\label{eq:MC_scalar_distension_deformation}
\end{equation}
%
Substitution of \eqref{eq:MC_scalar_distension_factors} into the material-stress
definition \eqref{eq:MC_solid_stress_definition} makes the isotropic
distension factors cancel.  The material stress therefore takes the standard
Cauchy-stress form
%
\begin{equation}
	\bs\sigma_s'
	=
	\rho_s
	\left.
	\p{\psi_s}{\bar{\mbf F}_s^e}
	\right\vert_{\mc X_s}
	(\bar{\mbf F}_s^e)^T,
	\qquad
	s\in\mc S .
	\label{eq:MC_scalar_distension_solid_stress}
\end{equation}
%
For a generic multicomponent solid, the tensor distension restriction
\eqref{eq:MC_distension_force_definition} reduces to the trace condition
%
\begin{equation}
	\frac{1}{3}\operatorname{tr}\left(
	\bs\sigma_s'
	\right)
	=
	\rho_s
	\left.
	\p{\psi_s}{a_s^e}
	\right\vert_{\mc X_s}
	a_s^e,
	\qquad
	s\in\mc S .
	\label{eq:MC_scalar_distension_stress_restriction}
\end{equation}
%
Using the phase-level Biot transform, the same condition can be written as a
restriction on the phase double-prime trace,
%
\begin{equation}
	\frac{1}{3}\operatorname{tr}(\bs\sigma_s^{\prime\prime})
	=
	\rho_s
	\left.
	\p{\psi_s}{a_s^e}
	\right\vert_{\mc X_s}
	a_s^e
	-\left(1-B_s\right)\bar p_E,
	\qquad
	s\in\mc S .
	\label{eq:MC_scalar_double_prime_distension_trace}
\end{equation}
%
These scalar relations are the reversible tensor restriction specialized to
\(a_s\).  The scalar specialization
also reduces the plastic-distension power in
\eqref{eq:MC_residual_dissipation} to
%
\begin{equation}
	\sum_{s\in\mc S}
	\bs\sigma_s':
	\mbf A_s^e\dot{\mbf A}_s^p
	(\mbf A_s^p)^{-1}(\mbf A_s^e)^{-1}
	=
	\sum_{s\in\mc S}
	\frac{1}{3}\operatorname{tr}\left(
	\bs\sigma_s'
	\right)
	\frac{\dot{a}_s^p}{a_s^p}.
	\label{eq:MC_scalar_plastic_distension_power}
\end{equation}
%
Thus scalar plastic distension is hydrostatic compaction or dilation driven by
the material mean stress \(\frac{1}{3}\operatorname{tr}(\bs\sigma_s')\).

\subsection{Distension and crystallization pressure}
\label{sec:MC_distension_crystallization_pressure}

Crystallization pressure enters
\eqref{eq:MC_overall_momentum_nonlinear_biot} through an admissible solid free
energy, an admissible kinetic law, or an algebraic affinity restriction.  In the
scalar setting, the volumetric split is
%
\begin{equation}
	\ln J=\ln a_s+\ln\bar J_s,
	\qquad
	s\in\mc S .
	\label{eq:MC_crystallization_volumetric_split}
\end{equation}
%
The scalar Coleman--Noll trace condition depends on the combined volumetric
strain \(\ln J=\ln a_s+\ln\bar J_s\).  The transformed fixed-pressure
volumetric potential therefore generates the double-prime stress from
\(\ln J\), with the Biot pressure correction supplied by
\eqref{eq:MC_phase_biot_split}.  A simple scalar choice, with tangent bulk
modulus \(K_s^\star\), is
%
\begin{equation}
	\widetilde W_{{\rm vol},s}
	=
	\frac{K_s^\star}{2}
	\left(\ln J\right)^2 .
	\label{eq:MC_component_transformed_volumetric_energy}
\end{equation}
%
Holding \(\bar J_s\) and \(\bar p_E\) fixed,
%
\begin{equation}
	a_s
	\left.
	\p{\widetilde W_{{\rm vol},s}}{a_s}
	\right\vert_{\bar J_s,\bar p_E}
	=
	K_s^\star\ln J .
	\label{eq:MC_component_transformed_volumetric_trace}
\end{equation}
%
Thus the transformed volumetric potential supplies
%
\begin{equation}
	\frac{1}{3}\operatorname{tr}(\bs\sigma_s^{\prime\prime})
	=
	K_s^\star\ln J ,
	\label{eq:MC_component_double_prime_trace_from_energy}
\end{equation}
%
and the original Helmholtz trace required by the scalar distension restriction is
%
\begin{equation}
	\rho_s
	\left.
	\p{\psi_s}{a_s^e}
	\right\vert_{\mc X_s}
	a_s^e
	=
	K_s^\star\ln J
	+\left(1-B_s\right)\bar p_E .
	\label{eq:MC_component_cn_trace_satisfied}
\end{equation}
%
This is \eqref{eq:MC_scalar_double_prime_distension_trace}.  The stress generated
by the energy supplies the crystallization-pressure contribution to momentum.

The mechanism affinity also supplies a prescribed or diagnostic
crystallization-pressure relation.  For a mechanism \((m)\) precipitating solid
phase \(s_{(m)}\),
%
\begin{equation}
	\mc A_{(m)}
	=
	-
	\sum_\xi\sum_\alpha
	\mu_\xi^\alpha\nu_{\xi (m)}^\alpha,
	\qquad
	\mc A_{(m)}\dot r_{(m)}\ge0 .
	\label{eq:MC_crystallization_affinity_kinetics}
\end{equation}
%
If a classical crystallization pressure \(\bar p_{\text{xtal},s_{(m)}}\) is used, it is
best interpreted as an algebraic affinity law,
%
\begin{equation}
	\mc A_{(m)}
	=
	\bar v_{s_{(m)}}
	\left[
	\bar p_{\text{xtal},s_{(m)}}
	-\frac{1}{3}
	\operatorname{tr}
	\left(
	\bs\sigma_{s_{(m)}}'
	\right)
	\right],
	\qquad
	\bar v_{s_{(m)}}=\frac{1}{\bar\rho_{s_{(m)}}}.
	\label{eq:MC_crystallization_affinity_volumetric}
\end{equation}
%
For example, the classical supersaturation formula
\(\bar p_{\text{xtal},s_{(m)}}=R\theta_{\mc F}\ln\Omega_{(m)}/V_{s_{(m)}}\),
where the fluid temperature is used because supersaturation and the
fluid--solid precipitation mechanism are evaluated in the mobile subsystem,
\(R\) is the gas constant, \(\Omega_{(m)}\) is the supersaturation ratio, and
\(V_{s_{(m)}}\) is the molar volume of the precipitating solid, gives
%
\begin{equation}
	\mc A_{(m)}
	=
	\bar v_{s_{(m)}}
	\left[
		\frac{R\theta_{\mc F}}{V_{s_{(m)}}}
		\ln\Omega_{(m)}
		-\frac{1}{3}
		\operatorname{tr}
		\left(
		\bs\sigma_{s_{(m)}}'
		\right)
		\right].
	\label{eq:MC_crystallization_supersaturation_affinity}
\end{equation}
%
At reversible equilibrium, \(\mc A_{(m)}=0\), and the classical
crystallization pressure equals the admissible current phase trace stress.
Solving the affinity and stress makes \(\Omega_{(m)}\) a diagnostic output.

\section{Multicomponent model summary}
\label{sec:MC_model_summary_closure_count}

This section records the unknown and equation counts for the complete model.
The tables describe a constrained mixed formulation in which all normalized
mass fractions, their multipliers, and both relative-flux families are retained
as fields.  Phase pressures and electrochemical potentials are evaluated from
the constitutive restrictions.  The solid phases share one skeleton
displacement, and the scalar distension factors are closed by
\eqref{eq:MC_solid_distension_evolution}.  The electroquasistatic formulation
includes the electric potential and Gauss' law.
Let \(P_{\mc F}=|\mc F|\), \(P_{\mc S}=|\mc S|\), and
\(P=P_{\mc F}+P_{\mc S}\).  Let \(N^{\mc F}\) be the number of tracked
components in each fluid phase and \(N^{\mc S}\) the number of tracked
components in each solid phase.  If the component counts vary by phase, the
products in the tables are replaced by the corresponding phase sums.  Let \(N_{(m)}\)
be the number of independent conversion mechanisms and \(d\) the spatial
dimension.

\begin{longtable}{>{\raggedright\arraybackslash}p{0.66\textwidth}>{\raggedleft\arraybackslash}p{0.24\textwidth}}
	\caption{Unknown families for a fully implicit pointwise discretization of the multicomponent model.}
	\label{tab:MC_unknown_family_count}                                                                                                                     \\
	Unknown family                                                                                  & Count                                                 \\
	\hline
	\endfirsthead
	\multicolumn{2}{l}{Table~\thetable\ continued}                                                                                                          \\
	Unknown family                                                                                  & Count                                                 \\
	\hline
	\endhead
	Fluid-phase displacements and the common solid-skeleton displacement                            & \(d(P_{\mc F}+1)\)                                    \\
	Intrinsic phase densities \(\bar\rho_\xi\) for \(\xi=1,\ldots,P\)                               & \(P\)                                                 \\
	Fluid mass fractions \(\eta_f^\alpha\)                                                          & \(P_{\mc F}N^{\mc F}\)                                \\
	Solid mass fractions \(\eta_s^\alpha\)                                                          & \(P_{\mc S}N^{\mc S}\)                                \\
	Phase volume fractions \(\phi_\xi\)                                                             & \(P\)                                                 \\
	Volume-fraction multiplier \(\lambda\)                                                          & \(1\)                                                 \\
	Composition-normalization multipliers \(\Pi_\xi\)                                               & \(P\)                                                 \\
	Storage multipliers \(\pi_\xi^\alpha\)                                                          & \(P_{\mc F}N^{\mc F}+P_{\mc S}N^{\mc S}\)             \\
	Transfer potential \(\tau\)                                                                     & \(1\)                                                 \\
	Dispersion and diffusion component-flux vectors                                                 & \(2d[P_{\mc F}(N^{\mc F}-1)+P_{\mc S}(N^{\mc S}-1)]\) \\
	Scalar solid distension factors \(a_s\), for \(s=1,\ldots,P_{\mc S}\)                           & \(P_{\mc S}\)                                         \\
	Mechanism process rates \(\dot r_{(1)},\ldots,\dot r_{(N_{(m)})}\)                              & \(N_{(m)}\)                                           \\
	Solid plastic internal variables \(\bar{\mbf F}_s^p\) and \(\mbf A_s^p\) for inelastic response & \(2d^2P_{\mc S}\)                                     \\
	Fluid and solid temperatures \(\theta_{\mc F}\), \(\theta_{\mc S}\)                             & \(2\)                                                 \\
	Electric potential \(\varphi\)                                                                  & \(1\)
\end{longtable}
%
The corresponding equation blocks are listed in
Table~\ref{tab:MC_equation_block_count}.
%
\begin{longtable}{>{\raggedright\arraybackslash}p{0.60\textwidth}>{\raggedleft\arraybackslash}p{0.30\textwidth}}
	\caption{Equation blocks and closure counts for the multicomponent model.}
	\label{tab:MC_equation_block_count}                                                                                                                                                             \\
	Equation block                                                                                                                          & Count                                                 \\
	\hline
	\endfirsthead
	\multicolumn{2}{l}{Table~\thetable\ continued}                                                                                                                                                  \\
	Equation block                                                                                                                          & Count                                                 \\
	\hline
	\endhead
	Fluid-phase momentum balances \eqref{eq:el_mom_f_simplified} and overall momentum balance \eqref{eq:MC_overall_momentum_nonlinear_biot} & \(d(P_{\mc F}+1)\)                                    \\
	Component mass balances \eqref{eq:averaged_component_spatial_mass}                                                                      & \(P_{\mc F}N^{\mc F}+P_{\mc S}N^{\mc S}\)             \\
	Dispersion and diffusion closures \eqref{eq:MC_relative_transport_closures}                                                             & \(2d[P_{\mc F}(N^{\mc F}-1)+P_{\mc S}(N^{\mc S}-1)]\) \\
	Component composition stationarity \eqref{eq:el_eta_all}                                                                                & \(P_{\mc F}N^{\mc F}+P_{\mc S}N^{\mc S}\)             \\
	Fluid and solid composition normalizations \eqref{eq:el_eta_constraint}                                                                 & \(P\)                                                 \\
	Volume-fraction restrictions \eqref{eq:MC_volume_fraction_restrictions} and constraint \eqref{eq:volume_fraction_constraint}            & \(P+1\)                                               \\
	Pressure-density restrictions \eqref{eq:MC_phase_pressure_definitions}                                                                  & \(P\)                                                 \\
	Evolution of \(\tau\) and component compatibility relations \eqref{eq:MC_admissible_conversion_component}                               & \(P_{\mc F}N^{\mc F}+P_{\mc S}N^{\mc S}\)             \\
	Solid distension evolution laws \eqref{eq:MC_solid_distension_evolution}                                                                & \(P_{\mc S}\)                                         \\
	Kinetic laws for \(\dot r_{(m)}\)                                                                                                       & \(N_{(m)}\)                                           \\
	Plastic flow laws for \(\bar{\mbf F}_s^p\) and \(\mbf A_s^p\) for inelastic response                                                    & \(2d^2P_{\mc S}\)                                     \\
	Fluid and solid energy equations using \eqref{eq:MC_energy_balance}                                                                     & \(2\)                                                 \\
	Gauss' law \eqref{eq:gauss_law}                                                                                                         & \(1\)
\end{longtable}
%
For a single-component solid phase, the normalization fixes
\(\eta_s^s=1\), and the component stationarity equation determines
\(\Pi_s\).  Equation \eqref{eq:MC_additive_component_biot} then gives
\(B_s^s=B_s\).  For one solid phase,
\eqref{eq:MC_solid_biot_rollup_coefficient} gives \(B=B_s\), and the
common-skeleton momentum equation is
\eqref{eq:MC_overall_momentum_nonlinear_biot}.
\FloatBarrier
